\documentclass[12pt]{book}
% Idioma y codificación
\usepackage[spanish, es-tabla]{babel}       %es-tabla para que se titule "Tabla"
\usepackage[utf8]{inputenc}

% Márgenes
\usepackage[a4paper,top=3cm,bottom=2.5cm,left=3cm,right=3cm]{geometry}

% Comentarios de bloque
\usepackage{verbatim}

% Paquetes de links
\usepackage[hidelinks]{hyperref}    % Permite enlaces
\usepackage{url}                    % redirecciona a la web

% Más opciones para enumeraciones
\usepackage{enumitem}

% Personalizar la portada
\usepackage{titling}


% Paquetes de tablas
\usepackage{multirow}


%------------------------------------------------------------------------

%Paquetes de figuras
\usepackage{caption}
\usepackage{subcaption} % Figuras al lado de otras
\usepackage{float}      % Poner figuras en el sitio indicado H.


% Paquetes de imágenes
\usepackage{graphicx}       % Paquete para añadir imágenes
\usepackage{transparent}    % Para manejar la opacidad de las figuras

% Paquete para usar colores
\usepackage[dvipsnames]{xcolor}
\usepackage{pagecolor}      % Para cambiar el color de la página

% Habilita tamaños de fuente mayores
\usepackage{fix-cm}


%------------------------------------------------------------------------

% Paquetes de matemáticas
\usepackage{mathtools, amsfonts, amssymb, mathrsfs}
\usepackage[makeroom]{cancel}     % Simplificar tachando
\usepackage{polynom}    % Divisiones y Ruffini
\usepackage{units} % Para poner fracciones diagonales con \nicefrac

\usepackage{pgfplots}   %Representar funciones
\pgfplotsset{compat=1.18}  % Versión 1.18

\usepackage{tikz-cd}    % Para usar diagramas de composiciones
\usetikzlibrary{calc}   % Para usar cálculo de coordenadas en tikz

%Definición de teoremas, etc.
\usepackage{amsthm}
%\swapnumbers   % Intercambia la posición del texto y de la numeración

\theoremstyle{plain}

\makeatletter
\@ifclassloaded{article}{
  \newtheorem{teo}{Teorema}[section]
}{
  \newtheorem{teo}{Teorema}[chapter]  % Se resetea en cada chapter
}
\makeatother

\newtheorem{coro}{Corolario}[teo]           % Se resetea en cada teorema
\newtheorem{prop}[teo]{Proposición}         % Usa el mismo contador que teorema
\newtheorem{lema}[teo]{Lema}                % Usa el mismo contador que teorema

\theoremstyle{remark}
\newtheorem*{observacion}{Observación}

\theoremstyle{definition}

\makeatletter
\@ifclassloaded{article}{
  \newtheorem{definicion}{Definición} [section]     % Se resetea en cada chapter
}{
  \newtheorem{definicion}{Definición} [chapter]     % Se resetea en cada chapter
}
\makeatother

\newtheorem*{notacion}{Notación}
\newtheorem*{ejemplo}{Ejemplo}
\newtheorem*{ejercicio*}{Ejercicio}             % No numerado
\newtheorem{ejercicio}{Ejercicio} [section]     % Se resetea en cada section


% Modificar el formato de la numeración del teorema "ejercicio"
\renewcommand{\theejercicio}{%
  \ifnum\value{section}=0 % Si no se ha iniciado ninguna sección
    \arabic{ejercicio}% Solo mostrar el número de ejercicio
  \else
    \thesection.\arabic{ejercicio}% Mostrar número de sección y número de ejercicio
  \fi
}


% \renewcommand\qedsymbol{$\blacksquare$}         % Cambiar símbolo QED
%------------------------------------------------------------------------

% Paquetes para encabezados
\usepackage{fancyhdr}
\pagestyle{fancy}
\fancyhf{}

\newcommand{\helv}{ % Modificación tamaño de letra
\fontfamily{}\fontsize{12}{12}\selectfont}
\setlength{\headheight}{15pt} % Amplía el tamaño del índice


%\usepackage{lastpage}   % Referenciar última pag   \pageref{LastPage}
\fancyfoot[C]{\thepage}

%------------------------------------------------------------------------

% Conseguir que no ponga "Capítulo 1". Sino solo "1."
\makeatletter
\@ifclassloaded{book}{
  \renewcommand{\chaptermark}[1]{\markboth{\thechapter.\ #1}{}} % En el encabezado
    
  \renewcommand{\@makechapterhead}[1]{%
  \vspace*{50\p@}%
  {\parindent \z@ \raggedright \normalfont
    \ifnum \c@secnumdepth >\m@ne
      \huge\bfseries \thechapter.\hspace{1em}\ignorespaces
    \fi
    \interlinepenalty\@M
    \Huge \bfseries #1\par\nobreak
    \vskip 40\p@
  }}
}
\makeatother

%------------------------------------------------------------------------
% Paquetes de cógido
\usepackage{minted}
\renewcommand\listingscaption{Código fuente}

\usepackage{fancyvrb}
% Personaliza el tamaño de los números de línea
\renewcommand{\theFancyVerbLine}{\small\arabic{FancyVerbLine}}

% Estilo para C++
\newminted{cpp}{
    frame=lines,
    framesep=2mm,
    baselinestretch=1.2,
    linenos,
    escapeinside=||
}



\usepackage{listings} % Para incluir código desde un archivo

\renewcommand\lstlistingname{Código Fuente}
\renewcommand\lstlistlistingname{Índice de Códigos Fuente}

% Definir colores
\definecolor{vscodepurple}{rgb}{0.5,0,0.5}
\definecolor{vscodeblue}{rgb}{0,0,0.8}
\definecolor{vscodegreen}{rgb}{0,0.5,0}
\definecolor{vscodegray}{rgb}{0.5,0.5,0.5}
\definecolor{vscodebackground}{rgb}{0.97,0.97,0.97}
\definecolor{vscodelightgray}{rgb}{0.9,0.9,0.9}

% Configuración para el estilo de C similar a VSCode
\lstdefinestyle{vscode_C}{
  backgroundcolor=\color{vscodebackground},
  commentstyle=\color{vscodegreen},
  keywordstyle=\color{vscodeblue},
  numberstyle=\tiny\color{vscodegray},
  stringstyle=\color{vscodepurple},
  basicstyle=\scriptsize\ttfamily,
  breakatwhitespace=false,
  breaklines=true,
  captionpos=b,
  keepspaces=true,
  numbers=left,
  numbersep=5pt,
  showspaces=false,
  showstringspaces=false,
  showtabs=false,
  tabsize=2,
  frame=tb,
  framerule=0pt,
  aboveskip=10pt,
  belowskip=10pt,
  xleftmargin=10pt,
  xrightmargin=10pt,
  framexleftmargin=10pt,
  framexrightmargin=10pt,
  framesep=0pt,
  rulecolor=\color{vscodelightgray},
  backgroundcolor=\color{vscodebackground},
}

%------------------------------------------------------------------------

% Comandos definidos
\newcommand{\bb}[1]{\mathbb{#1}}
\newcommand{\cc}[1]{\mathcal{#1}}

% I prefer the slanted \leq
\let\oldleq\leq % save them in case they're every wanted
\let\oldgeq\geq
\renewcommand{\leq}{\leqslant}
\renewcommand{\geq}{\geqslant}

% Si y solo si
\newcommand{\sii}{\iff}

% Letras griegas
\newcommand{\eps}{\epsilon}
\newcommand{\veps}{\varepsilon}
\newcommand{\lm}{\lambda}

\newcommand{\ol}{\overline}
\newcommand{\ul}{\underline}
\newcommand{\wt}{\widetilde}
\newcommand{\wh}{\widehat}

\let\oldvec\vec
\renewcommand{\vec}{\overrightarrow}

% Derivadas parciales
\newcommand{\del}[2]{\frac{\partial #1}{\partial #2}}
\newcommand{\Del}[3]{\frac{\partial^{#1} #2}{\partial^{#1} #3}}
\newcommand{\deld}[2]{\dfrac{\partial #1}{\partial #2}}
\newcommand{\Deld}[3]{\dfrac{\partial^{#1} #2}{\partial^{#1} #3}}


\newcommand{\AstIg}{\stackrel{(\ast)}{=}}
\newcommand{\Hop}{\stackrel{L'H\hat{o}pital}{=}}

\newcommand{\red}[1]{{\color{red}#1}} % Para integrales, destacar los cambios.

% Método de integración
\newcommand{\MetInt}[2]{
    \left[\begin{array}{c}
        #1 \\ #2
    \end{array}\right]
}

% Declarar aplicaciones
% 1. Nombre aplicación
% 2. Dominio
% 3. Codominio
% 4. Variable
% 5. Imagen de la variable
\newcommand{\Func}[5]{
    \begin{equation*}
        \begin{array}{rrll}
            #1:& #2 & \longrightarrow & #3\\
               & #4 & \longmapsto & #5
        \end{array}
    \end{equation*}
}

%------------------------------------------------------------------------

\theoremstyle{definition}
\newtheorem{cuestion}{Cuestión} [section]

% Para poder incluir árboles
\usepackage{forest}
\usepackage{booktabs}

% // TODO:
\newcommand{\showimages}{} % Descomentar esta línea para mostrar las imágenes

\begin{document}

    % 1. Foto de fondo
    % 2. Título
    % 3. Encabezado Izquierdo
    % 4. Color de fondo
    % 5. Coord x del titulo
    % 6. Coord y del titulo
    % 7. Fecha
    % 8. Autor

    
    % 1. Foto de fondo
% 2. Título
% 3. Encabezado Izquierdo
% 4. Color de fondo
% 5. Coord x del titulo
% 6. Coord y del titulo
% 7. Fecha

\newcommand{\portada}[7]{

    \portadaBase{#1}{#2}{#3}{#4}{#5}{#6}{#7}
    \portadaBook{#1}{#2}{#3}{#4}{#5}{#6}{#7}
}

\newcommand{\portadaExamen}[7]{

    \portadaBase{#1}{#2}{#3}{#4}{#5}{#6}{#7}
    \portadaArticle{#1}{#2}{#3}{#4}{#5}{#6}{#7}
}




\newcommand{\portadaBase}[7]{

    % Tiene la portada principal y la licencia Creative Commons
    
    % 1. Foto de fondo
    % 2. Título
    % 3. Encabezado Izquierdo
    % 4. Color de fondo
    % 5. Coord x del titulo
    % 6. Coord y del titulo
    % 7. Fecha
    
    
    \thispagestyle{empty}               % Sin encabezado ni pie de página
    \newgeometry{margin=0cm}        % Márgenes nulos para la primera página
    
    
    % Encabezado
    \fancyhead[L]{\helv #3}
    \fancyhead[R]{\helv \nouppercase{\leftmark}}
    
    
    \pagecolor{#4}        % Color de fondo para la portada
    
    \begin{figure}[p]
        \centering
        \transparent{0.3}           % Opacidad del 30% para la imagen
        
        \includegraphics[width=\paperwidth, keepaspectratio]{assets/#1}
    
        \begin{tikzpicture}[remember picture, overlay]
            \node[anchor=north west, text=white, opacity=1, font=\fontsize{60}{90}\selectfont\bfseries\sffamily, align=left] at (#5, #6) {#2};
            
            \node[anchor=south east, text=white, opacity=1, font=\fontsize{12}{18}\selectfont\sffamily, align=right] at (9.7, 3) {\textbf{\href{https://losdeldgiim.github.io/}{Los Del DGIIM}}};
            
            \node[anchor=south east, text=white, opacity=1, font=\fontsize{12}{15}\selectfont\sffamily, align=right] at (9.7, 1.8) {Doble Grado en Ingeniería Informática y Matemáticas\\Universidad de Granada};
        \end{tikzpicture}
    \end{figure}
    
    
    \restoregeometry        % Restaurar márgenes normales para las páginas subsiguientes
    \pagecolor{white}       % Restaurar el color de página
    
    
    \newpage
    \thispagestyle{empty}               % Sin encabezado ni pie de página
    \begin{tikzpicture}[remember picture, overlay]
        \node[anchor=south west, inner sep=3cm] at (current page.south west) {
            \begin{minipage}{0.5\paperwidth}
                \href{https://creativecommons.org/licenses/by-nc-nd/4.0/}{
                    \includegraphics[height=2cm]{assets/Licencia.png}
                }\vspace{1cm}\\
                Esta obra está bajo una
                \href{https://creativecommons.org/licenses/by-nc-nd/4.0/}{
                    Licencia Creative Commons Atribución-NoComercial-SinDerivadas 4.0 Internacional (CC BY-NC-ND 4.0).
                }\\
    
                Eres libre de compartir y redistribuir el contenido de esta obra en cualquier medio o formato, siempre y cuando des el crédito adecuado a los autores originales y no persigas fines comerciales. 
            \end{minipage}
        };
    \end{tikzpicture}
    
    
    
    % 1. Foto de fondo
    % 2. Título
    % 3. Encabezado Izquierdo
    % 4. Color de fondo
    % 5. Coord x del titulo
    % 6. Coord y del titulo
    % 7. Fecha


}


\newcommand{\portadaBook}[7]{

    % 1. Foto de fondo
    % 2. Título
    % 3. Encabezado Izquierdo
    % 4. Color de fondo
    % 5. Coord x del titulo
    % 6. Coord y del titulo
    % 7. Fecha

    % Personaliza el formato del título
    \pretitle{\begin{center}\bfseries\fontsize{42}{56}\selectfont}
    \posttitle{\par\end{center}\vspace{2em}}
    
    % Personaliza el formato del autor
    \preauthor{\begin{center}\Large}
    \postauthor{\par\end{center}\vfill}
    
    % Personaliza el formato de la fecha
    \predate{\begin{center}\huge}
    \postdate{\par\end{center}\vspace{2em}}
    
    \title{#2}
    \author{\href{https://losdeldgiim.github.io/}{Los Del DGIIM}}
    \date{Granada, #7}
    \maketitle
    
    \tableofcontents
}




\newcommand{\portadaArticle}[7]{

    % 1. Foto de fondo
    % 2. Título
    % 3. Encabezado Izquierdo
    % 4. Color de fondo
    % 5. Coord x del titulo
    % 6. Coord y del titulo
    % 7. Fecha

    % Personaliza el formato del título
    \pretitle{\begin{center}\bfseries\fontsize{42}{56}\selectfont}
    \posttitle{\par\end{center}\vspace{2em}}
    
    % Personaliza el formato del autor
    \preauthor{\begin{center}\Large}
    \postauthor{\par\end{center}\vspace{3em}}
    
    % Personaliza el formato de la fecha
    \predate{\begin{center}\huge}
    \postdate{\par\end{center}\vspace{5em}}
    
    \title{#2}
    \author{\href{https://losdeldgiim.github.io/}{Los Del DGIIM}}
    \date{Granada, #7}
    \thispagestyle{empty}               % Sin encabezado ni pie de página
    \maketitle
    \vfill
}
    \portada{ffccA4.jpg}{Ecuaciones\\Diferenciales I}{Ecuaciones Diferenciales I}{MidnightBlue}{-8}{28}{2024-2025}{José Juan Urrutia Milán}

    \newpage
    La parte de teoría del presente documento (es decir, excluyendo las relaciones de problemas) está hecha en función de los apuntes que se han tomado en clase.
    No obstante, recomendamos seguir de igual forma los apuntes del profesor de la asignatura, Rafael Ortega, disponibles en su sitio web personal
    \href{https://www.ugr.es/~rortega/}{https://www.ugr.es/~rortega/}. Estos apuntes no son por tanto una completa sustitución de dichos apuntes, sino tan solo un complemento.
    

    %\setcounter{chapter}{-1}
\chapter*{Introducción}

La teoría de Ecuaciones Diferenciales es la teoría matemática relacionada con el movimiento. Esta trata de resolver ecuaciones cuyas soluciones son funciones. Podríamos pensar en llamar a este área ``Ecuaciones Funcionales'', pero gracias a que en física $F = m \cdot a$, podemos encontrar de forma natural y útil ecuaciones funcionales donde la información que aparezca sobre la función a buscar está relacionada con las relaciones que guardan las derivadas de dicha función.

\begin{ejemplo}
    Como ejemplo de ecuación diferencial que surge de forma natural podemos pensar en el movimiento de un péndulo:

    Para describir un péndulo, nos es suficiente con tres variables independientes, que podemos ver en la siguiente ilustración:
    \begin{itemize}
        \item La longitud del hilo del péndulo, a la que llamamos $l$.
        \item La aceleración gravitatoria del planeta en el que nos encontremos, a la que llamamos $g$.
        \item Y el ángulo que guarda el péndulo sobre la vertical, al que llamamos $\theta$.
    \end{itemize}

    % Péndulo

\ifdefined\showimages
    \begin{figure}[H]
        \centering
        \begin{tikzpicture}
            % Parámetros
            \def\L{3}        % Longitud del péndulo
            \def\LL{2.7}     % Longitud de la vertical
            \def\ang{30}     % Ángulo del péndulo (en grados)
    
            % Dibujar la línea discontinua vertical
            \draw[dashed] (0,0) -- (0,-\LL);
    
            % Dibujar el péndulo
            \draw[thick] (0,0) -- ({\L*sin(\ang)}, {-\L*cos(\ang)});
    
            % Dibujar el círculo al final del péndulo
            \filldraw[thick] ({\L*sin(\ang)}, {-\L*cos(\ang)}) circle (0.2);
    
            % Indicar el ángulo
            \draw[thick, -stealth] (0,-1) arc[start angle=270, end angle=300, radius=1];
            \node at (0.35, -1.4) {$\theta$};
            \node at (1, -1) {$l$};
            \node at (3.3, -1) {$g$};

            % Flecha hacia abajo para la gravedad
            \draw[thick, -stealth] (3, -0.5) -- (3, -1.5);
    
            % Dibujar el punto de suspensión
            \filldraw (0,0) circle (0.05);
        \end{tikzpicture}
    \end{figure}
\fi
    Como nuestro objetivo es describir el movimiento que describe un péndulo, tenemos que introducir una variable más, el tiempo ($t$), y ver ahora la variable $\theta$ como una variable dependiente en función de $t$:
    \begin{equation*}
        \theta = \theta(t)
    \end{equation*}
    La física nos dice que si $\theta(t)$ es una función que nos describe el movimiento de un péndulo en función del tiempo, entonces debe cumplir la siguiente ecuación\footnote{De dónde sale dicha ecuación no nos es relevante.}:
    \begin{equation}\label{eq:pendulo}
        \theta''(t) + \frac{g}{l} \sen \theta(t) = 0
    \end{equation}

    A partir de esta ecuación, nos preguntamos por las funciones $\theta$ que cumplan dicha ecuación, siendo esta la primera ecuación diferencial que trataremos de resolver.

    A simple vista, podemos decir acertadamente que dos soluciones para dicha ecuación son:
    \begin{equation*}
        \theta(t) = 0 \qquad \qquad  \theta(t) = \pi
    \end{equation*}
    Pensando que en ambas soluciones el péndulo se encuentra en un estado estático, en la primera este se encuentra quieto debajo y en la segunda, quieto arriba. De forma intuitiva podemos pensar que el primero es un equilibrio estable y el segundo un equilibrio inestable.

    A partir de dichas soluciones, podemos adivinar que también serán soluciones de~(\ref{eq:pendulo}) cualquier función de la forma:
    \begin{equation*}
        \theta(t) = n\pi, \quad n \in \mathbb{Z}
    \end{equation*}
    De esta forma, hemos encontrado una \textbf{familia de soluciones}, es decir, tenemos una función que es solución de~(\ref{eq:pendulo}) para cada $n\in \mathbb{Z}$.

    Hemos encontrado infinitas soluciones para la ecuación~(\ref{eq:pendulo}). Podemos pensar que cada una de estas soluciones describe un péndulo distinto (esto es, cada una de las distintas formas de tirar el péndulo). En el mundo de las ecuaciones diferenciales es común encontrar muchas soluciones para una sola ecuación.

    En el caso de la ecuación~(\ref{eq:pendulo}), se ha demostrado que a parte de la familia de soluciones que hemos dado, no pueden encontrarse más soluciones con fórmula (aunque pueden aproximarse).
\end{ejemplo}~\\

El orden de una ecuación diferencial es la derivada de mayor grado qeu aparezca en la fórmula de la ecuación. En el caso de~(\ref{eq:pendulo}), esta era de orden 2. En esta asignatura nos centraremos en el estudio de las ecuaciones diferenciales de orden 1.

Una ecuación diferencial de primer orden genérica es de la forma:
\begin{equation}\label{eq_generica}
    \Phi(t,x(t), x'(t))=0
\end{equation}
Es decir, es una relación entre una variable independiente ($t$, que usualmente podremos entender como el tiempo), una variable dependiente o función ($x$, en función de $t$), cuya expresión estamos interesados en buscar; y su derivada.

La $\Phi$ que aparece en~(\ref{eq_generica}) será una función:
\Func{\Phi}{D\subseteq \mathbb{R}^3}{\mathbb{R}}{(t,x,y)}{\Phi(t,x,y)}
Donde trataremos de ver $x$ como variable independiente que tenemos que hacer dependiente de $t$: $x=x(t)$, siendo $y$ su derivada.

\begin{ejemplo}
    Dada la siguiente ecuación diferencial:
    \begin{equation*}
        {(x(t))}^{2} + {(x'(t))}^{2} = 1
    \end{equation*}
    La función $\Phi$ en cuestión es:
    \Func{\Phi}{D=\mathbb{R}^3}{\mathbb{R}}{(t,x,y)}{x^2+y^2-1}
    Soluciones de dicha ecuación son (a simple vista):

    \begin{equation*}
    \begin{array}{ll}
        x(t) = \sen t & x(t) = \cos t \\
        x(t) = 1 & x(t) = -1
    \end{array}
    \end{equation*}

    Además, podemos deducir que una familia de soluciones que engloba a las dos primeras es:
    \begin{equation*}
        x(t) = \sen(t+c), \quad c\in \mathbb{R}
    \end{equation*}

    % Grafica de familia de soluciones, 1 y -1
\ifdefined\showimages
    \begin{center}
    \begin{tikzpicture}
      \begin{axis}[
          axis lines = center,    % Dibujar los ejes
          xmin = -2*pi, xmax = 2*pi,   % Definir el rango en el eje x
          ymin = -1.5, ymax = 1.5,         % Definir el rango en el eje y
          xlabel = {},              % Etiqueta del eje x
          ylabel = {},                 % Etiqueta del eje y
          domain = -2*pi:2*pi,         % Dominio de las funciones
          samples = 200,               % Número de muestras para las funciones
          legend pos = north west,     % Posición de la leyenda
          %grid = both,                 % Mostrar la cuadrícula
          major grid style = {lightgray},   % Estilo de la cuadrícula principal
          minor grid style = {dotted}       % Estilo de la cuadrícula secundaria
      ]

      % Función constante 1
      \addplot[red, thick] {1};

      % Función constante -1
      \addplot[blue, thick] {-1};

      % Función sen(x)
      \addplot[green, thick] {sin(deg(x))};

      % Función cos(x)
      \addplot[purple, thick] {cos(deg(x))};

      % Función sen(x + pi)
      \addplot[purple, thick] {sin(deg(x + pi))};

      % Función sen(x + pi)
      \addplot[purple, thick] {sin(deg(x + 2 * pi))};

      \end{axis}
    \end{tikzpicture}
    \end{center}
\fi

    Graficando las soluciones, podemos además construir una nueva solución con una función a trozos:
    \begin{equation*}
        x(t) = \left\{\begin{array}{ll}
                \cos t & t \geq 0 \\
                1 & t<0
        \end{array}\right.
    \end{equation*}
    Derivable en $\mathbb{R}^*$ por el carácter local de la derivabilidad y en 0 por coincidir los dos límites laterales de las derivadas. Sin embargo, observamos que no es dos veces derivable.
\end{ejemplo}~\\

Las ecuaciones diferenciales de orden 1 sin restricción alguna son demasiado generales como para construir una teoría formal que centre su estudio en estas. Por tanto, estudiaremos aquellas que admitan escribirlas en \textbf{forma normal}, es decir, que si nos dan una ecuación en la forma~(\ref{eq_generica}), esta pueda escribirse como:
\begin{equation}\label{eq_generica_normal}
    x'(t) = f(t,x(t))
\end{equation}
Para una cierta
\Func{f}{C\subseteq \mathbb{R}^2}{\mathbb{R}}{(t,x)}{f(t,x)}

\begin{ejemplo}
    Dadas las dos siguientes ecuaciones diferenciales, discutir cuál de ellas es más simple y resolver ambas.
    \begin{itemize}
        \item $x'(t) = 7x(t)$
        \item $x'(t) = 7t$
    \end{itemize}
    La segunda ecuación es más sencilla, ya que se trata de un cálculo de una primitiva para $x'$: $x'(t) = h(t)$, a lo que ya estamos acostumbrados.

    Para dicha ecuación, tenemos:
    \begin{gather*}
        \Phi(t,x,y) = y-7t \\
        f(t,x) = 7t
    \end{gather*}

    con $D=\mathbb{R}^3$ y $C=\mathbb{R}^2$.

    Las soluciones de dicha ecuación diferencial son, por tanto:
    \begin{equation*}
        x(t) = \frac{7}{2} t^2 + c, \quad c\in \mathbb{R}
    \end{equation*}
    Para la primera ecuación, tenemos:
    \begin{gather*}
        \Phi(t,x,y) = y-7x \\
        f(t,x) = 7x
    \end{gather*}

    con $D=\mathbb{R}^3$ y $C=\mathbb{R}^2$.

    De forma simple vemos dos primeras soluciones:
    \begin{equation*}
        x(t) = 0 \qquad  x(t) = e^{7t}
    \end{equation*}

    Y podemos deducir además una familia de soluciones:
    \begin{equation*}
        x(t) = c \cdot e^{7t}, \quad c\in \mathbb{R}
    \end{equation*}

    De hecho, demostraremos más adelante que dicha ecuación no tiene más soluciones además de las de dicha familia.

\end{ejemplo}
   
\begin{ejemplo}
    Dadas las dos siguientes ecuaciones diferenciales, discutir cuál de ellas es más simple e intentar resolverlas.
    \begin{itemize}
        \item $x'(t) = \sen t$
        \item $x'(t) = \sen x(t)$
    \end{itemize}
    En este caso, es la primera la que es más simple, ya que se vuelve a tratar de un cálculo de primitiva.

    Tenemos:
    \begin{gather*}
        \Phi(t,x,y) = y-\sen t \\
        f(t,x) = \sen t
    \end{gather*}

    Siendo las soluciones de la ecuación diferencial:
    \begin{equation*}
        x(t) = -\cos t + c, \quad c \in \mathbb{R}
    \end{equation*}
    En el caso de la segunda, tenemos:
    \begin{gather*}
        \Phi(t,x,y) = y-\sen x \\
        f(t,x) = \sen x
    \end{gather*}
    Ecuación que todavía no sabemos resolver.
\end{ejemplo}

    \chapter{Ensamblador}
\section{Historia}
La linea de procesadores Intel, comúnmente conocida como \textit{x86}
ha seguido un largo proceso de evolución, manteniendo compatibilidad
hasta 8086, introducido en 1978. La familia de procesadores x86 ha seguido una evolución hasta
el día de hoy:
\begin{itemize}
	\item 8086 (1978) $\rightarrow$ 29K transistores, 16 bits, 5 MHz. \quad [Intel 16-bit, IBM PC]
	\item 386 (1985) $\rightarrow$ 275K transistores, 32 bits, 16 MHz. \quad [Intel 32-bit]
	\item Pentium 4E (2004) $\rightarrow$ 125M transistores, 64 bits, 3.4 GHz. \quad [Intel 64-bit, Hyperthreading, Implementación de x86-64 de AMD]
	\item Core i7 Haswell (2013) $\rightarrow$ 1.4G transistores, 64 bits, 3.9 Ghz \quad[Multicore, AVX2]
	\item Core i9 Raptor Lake (2022) $ \rightarrow$ +30G transistores, 64 bits, 5.2 GHz \quad[DDR5, PCIe 5.0, AVX-512]
\end{itemize}
A lo largo de los años, muchas compañias han creado procesadores
que se puedan comparar con los de Intel. Entre ellas, AMD ha sido
la más destacada, cayendo ligeramente detrás de Intel en cuanto a rendimiento
y ofreciendo un precio más competitivo. Sin embargo, en los últimos años, AMD ha
sido capaz de ofrecer procesadores con un rendimiento similar o superior a los de Intel,
concretamente, en 2024, Intel ha sufrido una gran decadencia por problemas de estabilidad
en los ultimos lanzamientos.
Nosotros nos centraremos en la arquitectura x86-64, que es la que se utiliza en el libro de
clase, sin embargo, haremos pequeñas menciones a IA32, cuando sea conveniente.
\section{Lenguaje C, ensamblador y máquina}
Las computadoras modernas utilizan distintas capas de abstracción, ocultando la complejidad
de la implementación. Dos de estas capas son realmente importantes para nosotros,
el formato y comportamiento de las instrucciones de máquina, que viene determinado
por la ISA (Instruction Set Architecture), esta define el formato de las instrucciones,
y el resultado que cada una de estas tendrá sobre el estado de la máquina.
La mayoria de ISAs, incluyendo x86-64, describen el comportamiento de un programa de
forma secuencial, como si se ejecutase una instrucción tras otra, sin embargo, en realidad,
el hardware de la CPU es mucho mas complejo, ejecutando instrucciones de forma paralelizada
o en otro orden, sin embargo, esto se verá con mas detenimiento en la asignatura \textit{Arquitectura de Computadores}.
La otra capa de abstracción ocurre sobre las direcciones de memoria, se nos presentan
como \textit{virtual addresses}, proporcionando un modelo de memoria que aparenta ser
un gran array de bytes, sin embargo, en realidad, la memoria se maneja de otra forma,
como se ve en la asignatura \textit{Sistemas Operativos}.
Recalcamos aunque no vayan a ser estudiadas en profundidad, que existen otras
ISA, como ARM (utilizada en dispositivos móviles), RISC-V (Nueva ISA open-source),
o la misma IA32. Aunque no se mencione más, se recomienda al lector la profundización
sobre ARM, que actualmente está ampliando su uso en el mercado de portatiles. (Intel Ultra, M1, etc.)
Recordamos que podemos clasificar en general las distintas ISAs en dos tipos:
\begin{itemize}
	\item \textbf{CISC} (Complex Instruction Set Computer): Conjunto de instrucciones complejas.
	\item \textbf{RISC} (Reduced Instruction Set Computer): Conjunto de instrucciones reducidas.
\end{itemize}
Realmente la mayoría de ISAs actuales son híbridas, convirtiendo
las instrucciones de CISC en instrucciones de RISC, para mejorar el rendimiento.
\newline
\newline
El codigo x86-64 difiere en gran medida del código C, destacamos las siguientes
partes del estado del procesador que normalmente no se ven en C:
\begin{itemize}
	\item \textbf{PC (Program Counter)}: Registra la dirección de la siguiente instrucción a ejecutar.
	      Se denomina \textit{\%rip} en x86-64.
	\item \textbf{Archivo de registros}: Contiene 16 registros etiquetados con valores de 64 bits.
	      Pueden ser utilizados para almacenar valores temporales, punteros, etc.
	\item \textbf{Codigos de condición}: Contienen información sobre la instrucción mas reciente.
	      Se utilizan para realizar saltos condicionales.
	\item \textbf{Pila}: Se utiliza para almacenar enteros o flotantes. Es un array direccionable por bytes.
\end{itemize}
Ahora, vamos a ver cómo se relaciona el código C con el código ensamblador. Para ello,
habrá que tener en cuenta que el código generado por el compilador puede variar
según el compilador, la versión, las optimizaciones, etc.
\subsection{Proceso de compilación}
Partimos de un número arbitrario de ficheros \textit{.c}, que contienen código C.
Para compilarlos, se pueden utilizar distintos compiladores, nosotros usaremos
\textit{gcc}:
% codigo en bash:
\begin{minted}{bash}
    gcc -Og programa.c -o programa
\end{minted}
Cuando lo ejecutemos:
\begin{center}
	\begin{tikzpicture}[node distance=1.6cm]

		% Styles within tikzpicture
		\tikzstyle{process} = [rectangle, rounded corners, minimum width=3cm, minimum height=1cm, text centered, draw=black, fill=green!30]
		\tikzstyle{arrow} = [thick,->,>=stealth]

		% Nodes
		\node (source) [process, fill=green!60] {Source File};
		\node (preprocessor) [process, below of=source] {Preprocessor};
		\node (compiler) [process, below of=preprocessor] {Compiler};
		\node (assembler) [process, below of=compiler] {Assembler};
		\node (linker) [process, below of=assembler] {Linker};
		\node (executable) [process, fill=green!60, below of=linker] {Executable};

		% Arrows
        \draw [-Stealth] (source) -- (preprocessor);
        \draw [-Stealth](preprocessor) -- node[right]{Preprocessed File (.i file)} (compiler);
        \draw [-Stealth](compiler) -- node[right]{Assembly code (.s file)} (assembler);
        \draw [-Stealth](assembler) -- node[right]{Object code (.o file)} (linker);
        \draw [-Stealth](linker) -- (executable);
	\end{tikzpicture}
\end{center}
En este proceso, gcc primero preprocesa el código, generando un archivo \textit{.i}, que
posteriormente compila, generando un archivo \textit{.s}, que se ensambla para generar
\textit{.o}, finalmente, se enlazan los distintos archivos \textit{.o} para generar el ejecutable.
Vamos a explicar un poco más detenidamente cada uno de los pasos:
\begin{itemize}
	\item \textbf{Preprocesador}: C proporciona ciertas herramientas del lenguaje mediante un preprocesador,
	      que es el primer paso en la compilación. Los más utilizados son \textit{\#include}, para incluir el contenido de un archivo
	      durante la compilación, y \textit{\#define}, para reemplazar un token por una secuencia arbitraria de texto
	      o valores. Además, elimina los comentarios de un archivo fuente. En gcc, podemos generar el código preprocesado con el flag \textit{-E}.
	\item \textbf{Compilación}: El compilador traduce el codigo preprocesado y genera código ensamblador, con gcc, podemos
	      generar el código ensamblador con el flag \textit{-S}.
	\item \textbf{Ensamblado}: El ensamblador traduce el código ensamblador a código de máquina, generando un archivo \textit{.o}.
	      En gcc, podemos generar el código de máquina con el flag \textit{-c}.
	\item \textbf{Enlazado}: El enlazador combina los distintos archivos \textit{.o} generados en un solo ejecutable. También se encarga de combinar
	      las librerias que se hayan utilizado en el código. En gcc, podemos generar el ejecutable con el flag \textit{-o}.
\end{itemize}


\begin{ejemplo}
	Veamos un ejemplo de cómo se relaciona el código C con el código ensamblador (solo se muestra la parte interesante):
	\begin{center} % Centramos el contenido
		\begin{minipage}{0.45\textwidth}
			\begin{minted}[fontsize=\Large]{c}
long plus(long x, long y);

void sumStore(long x, long y,
        long *dest) {
    long t = plus(x, y);
    *dest = t;
}
                \end{minted}
		\end{minipage}
		\begin{minipage}{0.45\textwidth}
			\begin{minted}[fontsize=\Large]{asm}
sumstore:
    pushq %rbx
    movq  %rdx, %rbx
    call  plus
    movq  %raw, (%rbx)
    popq  %rbx
    ret
                \end{minted}
		\end{minipage}
	\end{center}
\end{ejemplo}
Si lo compilamos en nuestra propia máquina, encontraremos muchas directivas que comienzan por . (ignorar)
\subsection{Representación Datos C, IA32, x86-64}
A continuación, mostramos los distintos tamaños en Bytes para los tipos de dato mas comúnes en C. Así como
su representación en Assembly. Debido a sus origenes como una arquitectura de 16 bits que expandió a 32 bits,
Intel usa el termino \textit{word} para referirse a un dato de 16 bits, y \textit{doubleword} para referirse a
una palabra de 32 bits. En x86-64, se ha añadido el termino \textit{quadword} para referirse a una palabra de 64 bits.
\begin{center}
	\begin{tabular}{|c|c|c|c|c|}
		\hline
		\textbf{Tipo de dato en C} & \textbf{Assembly sufijo} & \textbf{Normal 32 bits} & \textbf{IA32} & \textbf{x86-64} \\
		\hline
		\text{unsigned}            &                          & 4                       & 4             & 4               \\
		\text{int}                 & l                        & 4                       & 4             & 4               \\
		\text{long}                & q                        & 4                       & 4             & 8               \\
		\text{char}                & b                        & 1                       & 1             & 1               \\
		\text{short}               & w                        & 2                       & 2             & 2               \\
		\text{float}               & s                        & 4                       & 4             & 4               \\
		\text{double}              & l                        & 8                       & 8             & 8               \\
		\text{long double}         &                          & 8                       & 10/12         & 16              \\
		\text{char *}              &                          & 4                       & 4             & 8               \\
		\hline
	\end{tabular}
\end{center}
Los números de punto flotante se representan en el formato IEEE 754, ya visto en TOC. Se incluye
un diagrama de la representación de un número de punto flotante en simple y doble precision:
\begin{center}
	\includesvg[width=0.8\textwidth]{Imagenes/IEEE}
\end{center}
También cabe mencionar que existe un tipo de dato cuyo tamaño es 80 bits (long double) sin embargo su uso está
completamente desaconsejado debido a su portabilidad entre sistemas.
\section{Registros, operandos y operaciones}
\subsection{Registros}
Un procesador x86-64 tiene 16 registros de propósito general, que almacenan valores de 64 bits.
Estos registros se usan para almacenar valores temporales, punteros, etc. Sus nombres comienzan
con $\%r$, sin embargo, siguen convenciones de nombres debido a su legado. El primer 8086 tenía
8 registros de 16 bits, cada uno tenía un uso específico, de ahí su nombre, con la extensión a
IA32, los registros se extendieron a 32 bits, y con x86-64, a 64 bits, incluyendo además 8
nuevos registros, cuyos nombres van desde $\%r8$ hasta $\%r15$.

\begin{center}
	\includesvg[width=0.8\textwidth]{Imagenes/registers}
\end{center}
\subsection{Operandos}
La mayoría de instrucciones tienen uno o más operandos que especifican el \textit{source} y el destino de la operación.
En x86-64, los operandos pueden ser de tres tipos, (1) Inmediatos, (2) Registros, (3) Memoria. Veamos la sintaxis
para cada uno de ellos junto con la operación que se realiza:
\begin{center}
	\begin{tabular}{|l|l|l|l|}
		\hline
		\textbf{Type} & \textbf{Form}               & \textbf{Operand value}                      & \textbf{Name}       \\
		\hline
		Immediate     & \$\textit{Imm}              & \textit{Imm}                                & Immediate           \\
		Register      & $r_a$                       & $R[r_a]$                                    & Register            \\
		Memory        & \textit{Imm}                & $M[\textit{Imm}]$                           & Absolute            \\
		Memory        & $(r_a)$                     & $M[R[r_a]]$                                 & Indirect            \\
		Memory        & $\textit{Imm}(r_b)$         & $M[\textit{Imm} + R[r_b]]$                  & Base + displacement \\
		Memory        & $(r_b, r_i)$                & $M[R[r_b] + R[r_i]]$                        & Indexed             \\
		Memory        & $\textit{Imm}(r_b, r_i)$    & $M[\textit{Imm} + R[r_b] + R[r_i]]$         & Indexed             \\
		Memory        & $(r_i, s)$                  & $M[R[r_i] \cdot s]$                         & Scaled indexed      \\
		Memory        & $\textit{Imm}(, r_i, s)$    & $M[\textit{Imm} + R[r_i] \cdot s]$          & Scaled indexed      \\
		Memory        & $(r_b, r_i, s)$             & $M[R[r_b] + R[r_i] \cdot s]$                & Scaled indexed      \\
		Memory        & $\textit{Imm}(r_b, r_i, s)$ & $M[\textit{Imm} + R[r_b] + R[r_i] \cdot s]$ & Scaled indexed      \\
		\hline
	\end{tabular}
\end{center}
Los tipos de datos inmediatos se usan para valores constantes. En AT\&T se representan con un signo de dolar, siguiendo la notación
usual de C, por ejemplo \textit{\$0xFF} o \textit{\$323}. Los registros se representan con un signo de porcentaje, por ejemplo, $\%rax$.
Como ya hemos visto, los registros son de 64 bits, sin embargo, se pueden realizar operaciones de 32 bits, 16 bits o 8 bits, más adelante veremos cómo.
Los operandos de memoria se representan con paréntesis, por ejemplo, \textit{(\%rax)}, \textit{0x8(\%rax)}, etc. Podemos interpretarlo como una
desreferenciación de un puntero, es decir, acceder al valor almacenado en la dirección de memoria que se indica.
\begin{ejemplo}
	Veamos un ejemplo de cómo se realiza un swap en x86-64:
	\begin{center} % Centramos el contenido
		\begin{minipage}{0.45\textwidth}
			\begin{minted}[fontsize=\Large]{c}
void swap(long *xp, long *yp) 
{
    long x = *xp;
    *xp = *yp;
    *yp = x;
}
                \end{minted}
		\end{minipage}
		\begin{minipage}{0.45\textwidth}
			\begin{minted}[fontsize=\Large]{asm}
swap:
    movq (%rdi), %rax
    movq (%rsi), %rdx
    movq %rax, (%rsi)
    movq %rdx, (%rdi)
    ret
                \end{minted}
		\end{minipage}
	\end{center}
	Podemos ver que inicialmente, la disposición era $\% rdi \rightarrow xp$, $\% rsi \rightarrow yp$.
	Posteriormente, se almacena el \textbf{contenido} de $\%rdi$ en el registro $\%rax$, y el contenido de $\%rsi$ en $\%rdx$.
	Finalmente, se almacena el contenido de $\%rax$ en la dirección de memoria que apunta $\%rsi$, y el contenido de $\%rdx$ en la dirección de memoria que apunta $\%rdi$.
	Veamos un ejemplo para entender mejor el funcionamiento de los registros, supongamos que tenemos los siguientes valores:
	\begin{itemize}
		\item \textbf{Registro }$\%$rdi: 0x120
		\item \textbf{Registro }$\%$rsi: 0x130
		\item \textbf{Memoria en 0x120: } 125
		\item \textbf{Memoria en 0x130: } 450
	\end{itemize}
	Veamos ahora paso por paso el contenido de cada registro( $\rightarrow$ representa el contenido de la dirección de memoria):
	\[
		\begin{array}{|c|c|c|c|c|c|}
			\hline
			\textbf{Registro} & \textbf{Valor} & \textbf{movq}           & \textbf{movq}           & \textbf{movq}                                 & \textbf{movq}                         \\
			                  &                & \textbf{(\%rdi), \%rax} & \textbf{(\%rsi), \%rdx} & \textbf{\%rax, (\%rsi)}                       & \textbf{\%rdx, (\%rdi)}               \\
			\hline
			\%rdi             & 0x120          & 0x120                   & 0x120                   & 0x120                         \rightarrow 125 & 0x120 \textcolor{red}{\rightarrow450} \\
			\%rsi             & 0x130          & 0x130                   & 0x130                   & 0x130\textcolor{red}{\rightarrow125}          & 0x130 \rightarrow 125                 \\
			\%rax             & -              & \textcolor{red}{125}    & 125                     & 125                                           & 125                                   \\
			\%rdx             & -              & -                       & \textcolor{red}{450}    & 450                                           & 450                                   \\
			\hline
		\end{array}
	\]
	Es posible que el lector no entienda el por qué de las cuatro instrucciones en vez de tres como en C, esto se debe a que, como hemos visto antes, no se pueden realizar
	operaciones de memoria a memoria, por lo que se necesita un registro para almacenar el valor de una de las direcciones de memoria.
\end{ejemplo}
\subsection{Operaciones}
Veamos algunas de las operaciones más comunes en x86-64, notemos que algunas de ellas además pueden tener
variantes indicando el tamaño de los operandos. Antes de ver las operaciones de forma general, recalcamos
la instrucción \textbf{\textit{leaq}}:
\subsubsection{leaq}
La instrucción (\textit{load effective address}) se utiliza como variante de la instrucción \textit{movq},
la instrucción lee de la memoria la dirección efectiva de un operando y la almacena en un registro, sin acceder a la memoria en sí.
Por ejemplo \textit{leaq 0x8(\%rax), \%rdx} almacena en $\%rdx$ la dirección de memoria $\%rax + 0x8$. Esta instrucción es a menudo utilizada
por los compiladores para optimizaciones.
\begin{ejemplo}
	Veamos una forma para multiplicar un numero por 12 usando leaq
	\begin{center} % Centramos el contenido
		\begin{minipage}{0.3\textwidth}
			\begin{minted}[fontsize=\large]{c}
void m12(long x) 
{
    return x*12;
}
                \end{minted}
		\end{minipage}
		\begin{minipage}{0.69\textwidth}
			\begin{minted}[fontsize=\large]{asm}
swap:
    leaq (%rdi, %rdi, 2), %rax # t = x + 2x
    salq 2, %rax             # return t<<2
                \end{minted}
		\end{minipage}
	\end{center}
\end{ejemplo}

Veamos algunas de las operaciones del tipo aritmetico-lógicas
\begin{center}
    \begin{tabular}{|c|c|c|}
    \hline
    \textbf{Instrucción} & \textbf{Efecto}                       & \textbf{Descripción}     \\ \hline
    \texttt{lea S, D}         & $D \leftarrow \&S$                   & Load effective address   \\ \hline
    \texttt{inc D}            & $D \leftarrow D + 1$                 & Increment                \\ \hline
    \texttt{dec D}            & $D \leftarrow D - 1$                 & Decrement                \\ \hline
    \texttt{neg D}            & $D \leftarrow -D$                    & Negate                   \\ \hline
    \texttt{not D}            & $D \leftarrow \sim D$                & Complement               \\ \hline
    \texttt{add S, D}         & $D \leftarrow D + S$                 & Add                      \\ \hline
    \texttt{sub S, D}         & $D \leftarrow D - S$                 & Subtract                 \\ \hline
    \texttt{imul S, D}        & $D \leftarrow D * S$                 & Multiply                 \\ \hline
    \texttt{xor S, D}         & $D \leftarrow D \oplus S$            & Exclusive-or             \\ \hline
    \texttt{or S, D}          & $D \leftarrow D \vee S$              & Or                       \\ \hline
    \texttt{and S, D}         & $D \leftarrow D \wedge S$            & And                      \\ \hline
    \texttt{sal k, D}         & $D \leftarrow D \ll k$               & Left shift               \\ \hline
    \texttt{shl k, D}         & $D \leftarrow D \ll k$               & Left shift (same as sal) \\ \hline
    \texttt{sar k, D}         & $D \leftarrow D \gg_A k$             & Arithmetic right shift   \\ \hline
    \texttt{shr k, D}         & $D \leftarrow D \gg_L k$             & Logical right shift      \\ \hline
    \end{tabular}
\end{center}

Destacamos una característica peculiar de las operaciones de shift, en el caso de que hagamos un shift a la derecha sobre un
registro de 32 bits \texttt{shr \$16 $\%$edx}, $\%$rdx se rellenará con 0 en los primeros 32 bits también, de esta forma si teniamos $\$-1$
por ejemplo y realizamos \texttt{shr \$16 $\%$edx}, nos quedaremos con 0x0000 0000 0000 FFFF en vez de 0xFFFF FFFF 0000 FFFF.
\subsubsection{Operaciones aritméticas especiales}
Multiplicar dos enteros con/sin signo de 64 bits puede tener un resultado que necesite 128 bits para representarlo
es para ello que se incluyen las siguientes instrucciones especiales:
\begin{center}
    \begin{tabular}{|c|c|c|}
    \hline
        \textbf{Instruction} & \textbf{Effect} & \textbf{Description} \\ \hline
    \texttt{imulq S} & \texttt{R[\%rdx]:R[\%rax] ← S × R[\%rax]} & Signed full multiply \\ \hline
    \texttt{mulq S} & \texttt{R[\%rdx]:R[\%rax] ← S × R[\%rax]} & Unsigned full multiply \\ \hline
    \texttt{cqto} & \texttt{R[\%rdx]:R[\%rax] ← SignExtend(R[\%rax])} & Convert to oct word \\ \hline
    \texttt{idivq S} & 
    \begin{tabular}[t]{@{}l@{}}
        \texttt{R[\%rdx] ← R[\%rdx]:R[\%rax] mod S;} \\ 
        \texttt{R[\%rax] ← R[\%rdx]:R[\%rax] ÷ S}
    \end{tabular} & Signed divide \\ \hline
    \texttt{divq S} & 
    \begin{tabular}[t]{@{}l@{}}
        \texttt{R[\%rdx] ← R[\%rdx]:R[\%rax] mod S;} \\ 
        \texttt{R[\%rax] ← R[\%rdx]:R[\%rax] ÷ S}
    \end{tabular} & Unsigned divide \\ \hline
    \end{tabular}
\end{center}
Además de esto, añadimos ciertas peculiaridades que nos pueden ser útiles en la práctica:

\begin{itemize}
    \item Las instrucciones \texttt{movzbl} y \texttt{movslq} (la \texttt{z} indica \textit{zero extension} y la \texttt{s} \textit{sign extension}) se utilizan para extender cantidades a registros. Lo remarcable es que, mientras que \texttt{movslq} mueve 32 bits a 64 bits con extensión de signo, \texttt{movzlq} no existe ya que al mover 32 bits a 64 bits (\texttt{movl}) ya se realiza una extensión de ceros.
    \item Cuando se tiene el código desensamblado de un programa, podemos ver la dirección en hexadecimal, junto con varios pares de números en hexadecimal y a la derecha las instrucciones. Los pares de números en hexadecimal representan los bytes de la instrucción. Por ejemplo, si tenemos la línea siguiente:
    \begin{center}
        400544: \quad \texttt{e8 08 00 00 00} \quad \texttt{callq 400551 <f>}
    \end{center}
    Podremos ver que \texttt{e8} corresponde con la instrucción \texttt{callq} y los siguientes 4 bytes son la dirección a la que se llama. En este caso, \texttt{400544} + 8 = \texttt{400551}, como se puede ver en la instrucción a la derecha.
\end{itemize}


\section{Control}
Hasta ahora hemos considerado como se comporta el código cuando se ejecuta de forma secuencial, sin embargo,
algunas estructuras de $C$ como los bucles, las estructuras de control, etc. requieren de una ejecución condicional, esto es,
la secuencia de instrucciones a ejecutar depende de una condición.
En x86-64, veremos dos formas de implementar estas estructuras de control más adelante, primero, veamos
los códigos de condición que se utilizan en x86-64:
\subsection{Codigos de condición}
Además de los registros de propósito general, la CPU mantiene un conjunto de bits de estado que se utilizan para
describir atributos de la última operación aritmético-lógica realizada.
Veamos los códigos de condición que vamos a usar:
\begin{itemize}
	\item \textbf{CF (Carry Flag)}: El flag de acarreo, se activa si hay un acarreo en la última operación. Se utiliza para detectar overflow en operaciones sin signo.
	\item \textbf{ZF (Zero Flag)}: El flag de cero, se activa si el resultado de la última operación es cero.
	\item \textbf{SF (Sign Flag)}: El flag de signo, se activa si el resultado de la última operación es negativo.
	\item \textbf{OF (Overflow Flag)}: El flag de desbordamiento, se activa si hay un desbordamiento en la última operación en complemento a 2.
\end{itemize}
Estos códigos de condición se ajustan implicitamente por las operaciones aritméticas, de esta forma, \textit{leaq} no afecta a los códigos de condición.
Además de ajustarse por las operaciones que ya conocemos, existen instrucciones que ajustan explicitamente los códigos de condición sin alterar los registros de propósito general.
\begin{itemize}
	\item El conjunto de instrucciones \textbf{\textit{Compare}} (cmp) ajustan los códigos de condición de acuerdo a las diferencias entre sus dos operandos. Debemos tener en cuenta que en $AT\&T$ la sintaxis
	      es al revés: \textit{cmpq S1, S2} compara $S1$ con $S2$ mediante $S2- S1$.
	\item El conjunto de instrucciones \textbf{\textit{Test}} (test) ajustan los códigos de condición de acuerdo a las operaciones lógicas entre sus dos operandos.
	      Se realiza la operación $S1 \& S2$ y se ajustan los códigos de condición. Es común que se utilicen ambos operandos iguales, por ejemplo, para
	      comprobar si un número es 0.
\end{itemize}
\subsubsection{Consultar códigos de condición}
En vez de consultarlos directamente, se utilizan mediante tres formas:
\begin{enumerate}
	\item Se establece un único byte a 0 o 1 en función de una combinación de los códigos de condición.
	\item Se utilizan para realizar saltos condicionales.
	\item Se utilizan para transferir datos condicionalmente.
\end{enumerate}
Para la primera forma, se utilizan las instrucciones \textit{setcc} siguientes:
\begin{center}
	\begin{tabular}{|c|c|}
		\hline
		\textbf{Instruction} & \textbf{Effect}                                                          \\ \hline
		\texttt{sete D}      & D $\leftarrow$ ZF (Equal / zero)                                         \\ \hline
		\texttt{setne D}     & D $\leftarrow$ $\sim$ ZF (Not equal / not zero)                          \\ \hline
		\texttt{sets D}      & D $\leftarrow$ SF (Negative)                                             \\ \hline
		\texttt{setns D}     & D $\leftarrow$ $\sim$ SF (Nonnegative)                                   \\ \hline
		\texttt{setg D}      & D $\leftarrow$ $\sim$ (SF $\oplus$ OF) \& $\sim$ ZF (Greater, signed >)  \\ \hline
		\texttt{setge D}     & D $\leftarrow$ $\sim$ (SF $\oplus$ OF) (Greater or equal, signed $\geq$) \\ \hline
		\texttt{setl D}      & D $\leftarrow$ SF $\oplus$ OF (Less, signed <)                           \\ \hline
		\texttt{setle D}     & D $\leftarrow$ (SF $\oplus$ OF) $\lor$ ZF (Less or equal, signed $\leq$) \\ \hline
		\texttt{seta D}      & D $\leftarrow$ $\sim$ CF \& $\sim$ ZF (Above, unsigned >)                \\ \hline
		\texttt{setae D}     & D $\leftarrow$ $\sim$ CF (Above or equal, unsigned $\geq$)               \\ \hline
		\texttt{setb D}      & D $\leftarrow$ CF (Below, unsigned <)                                    \\ \hline
		\texttt{setbe D}     & D $\leftarrow$ CF $\lor$ ZF (Below or equal, unsigned $\leq$)            \\ \hline
	\end{tabular}
\end{center}
\begin{ejemplo}
	Veamos un ejemplo de su uso:
	\begin{center}
		\begin{minipage}{0.39\textwidth}
			\begin{minted}[fontsize=\Large]{c}
    long gt(long a, long b)
    {
        return a > b;
    }
            \end{minted}
		\end{minipage}
		\begin{minipage}{0.6\textwidth}
			\begin{minted}[fontsize=\Large]{asm}
    gt:
        cmpq %rsi, %rdi
        setg %al
        movzbq %al, %rax
        ret

            \end{minted}
		\end{minipage}
	\end{center}
	Donde movzbq es una instrucción que significa \textit{ Move with Zero-extend Byte to Long}.
\end{ejemplo}

Veamos ahora cómo se realizan los saltos condicionales.
Estos provocan que la ejecución del programa pueda saltar a otra parte del código en función de los códigos de condición, es muy posible que
en algún momento nos encontremos una instrucción en C de tipo goto, su uso está desaconsejado, sin embargo, por su similitud con el ensamblador,
puede resultar útil en ciertas situaciones. Veamos una tabla con las instrucciones de salto condicional y después un ejemplo de uso.

\begin{center}
	\begin{tabular}{|c|c|c|}
		\hline
		\textbf{Instruction} & \textbf{Jump Condition}                 & \textbf{Description}         \\ \hline
		jmp                  & Label                                   & Direct jump                  \\ \hline
		jmp                  & *Operand                                & Indirect jump                \\ \hline
		je                   & ZF                                      & Equal / zero                 \\ \hline
		jne                  & $\sim$ZF                                & Not equal / not zero         \\ \hline
		js                   & SF                                      & Negative                     \\ \hline
		jns                  & $\sim$SF                                & Nonnegative                  \\ \hline
        jg                   & $\sim$(SF \^{} OF) $\&$ $\sim$ZF & Greater (signed >)           \\ \hline
        jge                  & $\sim$(SF \^{} OF)                  & Greater or equal (signed >=) \\ \hline
        jl                   & SF \^{} OF                          & Less (signed <)              \\ \hline
        jle                  & (SF \^{}  OF) $|$ ZF              & Less or equal (signed <=)    \\ \hline
        ja                   & $\sim$CF $\&$ $\sim$ZF               & Above (unsigned >)           \\ \hline
		jae                  & $\sim$CF                                & Above or equal (unsigned >=) \\ \hline
		jb                   & CF                                      & Below (unsigned <)           \\ \hline
		jbe                  & CF $|$ ZF                            & Below or equal (unsigned <=) \\ \hline
	\end{tabular}
\end{center}
\begin{ejemplo}
	Veamos ahora un ejemplo de cómo hacer por ejemplo un valor absoluto de una diferencia:
	\begin{center}
		\begin{minipage}{0.39\textwidth}
			\begin{minted}[fontsize=\Large]{c}
    long abs(long x, long y)
    {
        if (x > y)
            return x - y;
        else
            return y - x;
    }
            \end{minted}
		\end{minipage}
		\begin{minipage}{0.6\textwidth}
			\begin{minted}[fontsize=\Large]{asm}
        abs:
            cmpq %rsi, %rdi
            jle .L1
            movq    %rdi, %rax
            subq    %rsi, %rax
            ret
        .L1:
            movq %rsi, %rax
            subq %rdi, %rax
            ret
			\end{minted}
		\end{minipage}
	\end{center}
	Si tratamos de \href{https://godbolt.org/z/aj7f7os8d}{replicar} el ejemplo anterior, veremos que en general, el compilador no produce este código,
	esto se debe a las optimizaciones que el compilador aplica, evitando en la medida de lo posible los saltos condicionales, ya que muchas veces,
	es mejor calcular ambas ramas, a tener que predecir y calcular cual de las ramas hay que seguir.
	Para obtener este mismo código, deberemos especificar la opción de compilación \textit{-fno-if-conversion}.
\end{ejemplo}
Nos detenemos brevemente en las optimizaciones mencionadas, lo que hace el compilador en vez de saltar ramas es asignar valores a las variables
en función de un test: \texttt{let val = test ? x : y}. \textit{gcc} intenta utilizar este tipo de optimizaciones siempre que sea seguro.
De esta forma, el código anterior se traduciría a:
\begin{center}
	\begin{minted}[fontsize=\Large]{asm}
        abs:
            movq %rdi, %rax # x
            subq %rsi, %rax # x - y
            movq %rsi, %rdx 
            subq %rdi, %rdx # y - x
            cmpq %rsi, %rdi # x : y
            cmovle %rdx, %rax # si x <= y, rax = rdx
    \end{minted}
\end{center}
Esto se verá como una optimización en la asignatura de \textit{Arquitectura de Computadores}, sin embargo, adelantamos
que por ejemplo \texttt{ let val = test ? p : *q} no se puede optimizar de la misma forma, ya que pudiera tener efectos no deseables.
Notemos además que este tipo de optimizaciones solo tienen sentido cuando los cálculos son sencillos, en cuyo caso es preferible calcular las
dos ramas a predecir cual de ellas seguir. Es por este tipo de opciones que \textit{gcc} es posible que genere un código mas lento cuando se activan
las optimizaciones \textit{-O3} ya que si la rama era muy predecible, la predicción de saltos es mas rápida que calcular las dos ramas.
Tampoco se debe utilizar en casos donde una rama u otra tiene efectos colaterales, ya que eso pudiera cambiar el comportamiento del programa.
\subsection{Bucles}
La mayoría de lenguajes ofrecen varias formas de implementar bucles: \textit{ do-while, while, for, loop, etc.} En x86-64, no existen instrucciones
similares, en su lugar, se combinan diferentes tests y saltos condicionales para implementar los bucles.
\subsubsection{Bucles Do-While}
En C, un bucle \textit{do-while} se implementa de la siguiente forma:
\begin{center}
	\begin{minipage}{0.45\textwidth}
		\begin{minted}[fontsize=\Large]{c}
long pcount(long x)
{
    long result = 0;
    do
    {
        result += x & 0x1;
        x >>= 1;
    } while (x);
    return result;
}
        \end{minted}
		\descriptionlabel{Versión C}
	\end{minipage}
	\begin{minipage}{0.45\textwidth}
		\begin{minted}[fontsize=\Large]{c}
long pcount(long x)
{
    long result = 0;
    loop:
        result += x & 0x1;
        x >>= 1;
        if(x) goto loop;
    return result;
}
        \end{minted}
		\descriptionlabel{Versión goto}

	\end{minipage}
\end{center}
En este ejemplo se ha hecho un clásico programa popcount que cuenta el número de bits a 1 en un número.
El nombre \textit{popcount} viene de \textit{population count}, que es el nombre que se le da a esta operación.
Como ya sabemos, el efecto del loop es ejecutar el cuerpo del bucle, evaluar la condición y si el resultado de la condición es distinto de 0,
seguir ejecutando el bucle. Notemos que se ejecuta al menos una vez.
\subsubsection{Bucles While}
\begin{ejemplo}
	Veamos ahora cómo se implementa un bucle \textit{while}:
	\begin{center}
		\begin{minipage}{0.45\textwidth}
			\begin{minted}[fontsize=\Large]{c}
long pcount(long x)
{
    long result = 0;
    while(x) {
        result += x & 0x1;
        x >>= 1;
    }

    return result;
}
\end{minted}
			\descriptionlabel{Versión C}
		\end{minipage}
		\begin{minipage}{0.45\textwidth}
			\begin{minted}[fontsize=\Large]{c}

long pcount(long x)
{
        long result = 0;
        goto test;
    loop:
        result += x & 0x1;
        x >>= 1;
    test:
        if (x) goto loop;
        return result;
}
        \end{minted}
			\descriptionlabel{Versión goto}

		\end{minipage}
	\end{center}
\end{ejemplo}
En general, while evalúa la condición antes de ejecutar el bucle, por lo que si la condición no se cumple, el bucle no se ejecuta.
\subsubsection{Bucles For}
Sabemos ya que en general un bucle \textit{for} se compone de tres partes, la inicialización, la condición y la actualización.
No adjuntamos el código en C ya que es trivial sabiendo los anteriores, notemos que un bucle \textit{for} es equivalente a un bucle \textit{while} de la siguiente forma:
\textit{for (init; test; update) body} es equivalente a \textit{init; while(test) {body; update;}}
Por último, veamos las secuencias switch.
\subsection{Secuencias Switch}
Una secuencia switch proporciona una ramificación condicional múltiple basada en el valor de una expresión entera.
Son muy útiles cuando trabajamos con tests que pueden tener una gran cantidad de resultados, haciendo el código más legible
y eficiente. Este último punto puede sorprender, por ello, vamos a desarrollarlo.
\\ \\
Una \textit{tabla de salto} es una estructura de datos en forma de array donde $i$ es la dirección de un segmento de codigo que ejecuta la acción
que el programa debería seguir cuando el valor de la expresión es $i$. La diferencia entre multiples \textit{if-else} y \textit{switch} radica en
que el tiempo que se tarda en ejecutar un \textit{switch} es constante, mientras que en un \textit{if-else} es lineal en función del número de casos.

\begin{center}
	\begin{minipage}{0.45\textwidth}
		\centering
		\textbf{Código en C}

		\begin{minted}[fontsize=\large]{c}

void switch_eg(long x) {
    switch (x) {
        case 1:
            printf("Code 1\n");
            break;
        case 2:
            printf("Code 2\n");
            break;
        case 3:
            printf("Code 3\n");
            break;
        default:
            printf("Code 4\n");
    }
}

    \end{minted}

	\end{minipage}
	\hfill
	\begin{minipage}{0.45\textwidth}
		\centering
		\textbf{Assembly Code}

		\begin{minted}[fontsize=\large]{asm}
cmpq $3, %rdi
ja .L4 # Si x > 3, Default
jmp *.L1(,%rdi,8) # Dirección L1 + 8x 
.L1:
.L2:
.L3:
.L4 # Default
        \end{minted}

		\vspace{1em}
		\textbf{Branch Table (Jump Table)}

		\begin{tabular}{|c|c|}
			\hline
			\textbf{Índice} & \textbf{Dirección de Salto} \\
			\hline
			1               & Code Block 1                \\
			\hline
			2               & Code Block 2                \\
			\hline
			3               & Code Block 3                \\
			\hline
			Default         & Code Block 4                \\
			\hline
		\end{tabular}
	\end{minipage}
\end{center}
Como podemos ver, la tabla de saltos comienza en este caso en .L1, para acceder
a el resto de bloques de código, se suma 8 veces el valor de x, ya que cada dirección
ocupa 8 bytes. Si x es mayor que 3, se salta a la dirección .L4, que es el bloque por defecto.
Si por ejemplo en vez de hacer break, se quisiera hacer un fall through, obtendríamos algo
del siguiente estilo:
\begin{center}
	\begin{minipage}{0.45\textwidth}
		\begin{minted}[fontsize=\large]{c}
long w = 1;
switch(x) {
    ...
  case 2:
    w = y/z;
    /* Fall through */ 
case 3:
    w += z;
    break;
    ...
}
\end{minted}
	\end{minipage}
	\begin{minipage}{0.45\textwidth}
		\begin{equation*}
			\%rdi \rightarrow x \quad \%rsi \rightarrow y \quad \%rdx \rightarrow z
		\end{equation*}
		\begin{minted}[fontsize=\large]{asm}
.L5: # Case 2
    movq %rsi, %rax
    cqto # Convert Quad To Oct
    idivq %rdx 
    jmp .L6
.L9: # Case 3
    movl $1, %eax
.L6: # Merge
    addq %rcx, %rax # w += z
    ret 

\end{minted}
	\end{minipage}
\end{center}

\section{Procedimientos}
Los procedimientos son una abstracción clave en el desarrollo de software, permiten dividir un programa en pequeñas partes
donde cada una de ellas proporciona una funcionalidad concreta al proporcionarle una serie de argumentos.
Pueden ser llamados desde cualquier parte del programa, lo que permite reutilizar el código y hacerlo más modular. Los procedimientos vienen
de distintas formas en cada lenguaje: funciones, métodos, subrutinas, lambdas, etc, pero todos ellos comparten ciertos elementos comunes.
\\ \\
Existen muchos atributos que hay que manejar cuando se proporciona soporte en lenguaje ensamblador para procedimientos, para verlos, supongamos un proceso
$P$ que llama a un procesor $Q$, que se ejecuta y tras esto vuelve a $P$, estas acciones incluyen:
\begin{itemize}
	\item \textbf{Pasar Control}: El PC tiene que fijarse en la dirección de comienzo de $Q$, y tras esto, cuando $Q$ termina, el PC tiene que volver a la dirección de la instrucción que sigue a la llamada.
	\item \textbf{Parámetros}: $P$ tiene que pasar los parámetros a $Q$, y $Q$ tiene que devolver un valor a $P$.
	\item \textbf{Reserva y Limpieza de Memoria}: $Q$ puede necesitar reservar memoria para sus propias variables, y tras la ejecución, limpiarla.
\end{itemize}
La implementación de estos procedimientos en x86-64 involucra una combinación de instrucciones especiales y convenciones de cómo usar los recursos de la máquina.
Detrás de esto hay un gran esfuerzo por minimizar el overhead que involucra una llamada, esto provoca que al final, solo se ejecuta lo mínimo indispensable mencionado
antes. Veamos ahora diferentes mecanismos paso por paso, comenzando por la pila:
\subsection{Stack (Pila)}
Es una región de memoría gestionada con aspectos de LIFO, es decir, el último elemento en entrar es el primero en salir, crece hacia direcciones inferiores, en x86-64
el registro $\%rsp$ contiene la dirección mas baja de la pila. Los datos, pueden ser accedidos y extraidos de la pila usando las instrucciones \textit{pushq} y \textit{popq}.
Por ejemplo para reservar espacio para datos sin valor inicial, se puede decrementar la pila la cantidad necesaria, análogamente, se puede liberar el espacio incrementando el stack pointer.
En x86-64, cuando un procedimiento requiere más espacio de lo que puede guardar en registros, reserva espacio en el stack. A esta región se la llama \textit{procedure space} para las variables locales.
En interés del espacio y la eficiencia, los procedimientos de x86-64 reservan solo el espacio necesario.

\subsection{Pasando el Control}
Pasar el control de una función $P$ a otra $Q$ requiere únicamente de poner el PC en la dirección de inicio de $Q$,
sin embargo, cuando $Q$ termina, el PC tiene que tener alguna manera de saber dónde debe reanudar la ejecución de $P$, esta información se guarda en las máquinas
x86-64 usando la instrucción call $Q$. Esta instrucción añade una dirección $A$ al stack y pone el PC en el comienzo de $Q$. La dirección $A$ se denomina \textit{return address} y se calcula como
la instrucción que sigue inmediatamente a call. La instrucción contraria ret, elimina una dirección $A$ de la pila y pone el PC en $A$.
(Notemos que estas instrucciones nos las encontraremos como callq y retq, donde el sufijo $q$ simplemente enfatiza que son versiones de x86-64 no de IA32).
Veamos un esquema con la estructura general del stack:


\begin{center}
	\includesvg[width=0.6\textwidth]{Imagenes/stack}
	\label {fig:stack}
\end{center}
\subsection{Pasando los datos}
Además de pasar el control a un procedimiento cuando es invocado, y de vuelta cuando termina, las llamadas a procedimientos
pueden involucrar pasar datos como argumentos, además un procedimiento puede devolver también un valor. En x86-64, la mayoria de
este flujo de datos es llevada a cabo mediante registros, ya hemos visto numerosos ejemplos en los que $\%rdi$ y $\%rsi$ se pasan como
parámetros y valores que se devuelven mediante $\%rax$. En x86-64, se pueden pasar hasta 6 argumentos mediante registros, que se usan en un orden
específico con el nombre del registro en función del tamaño del dato que contiene.
Cuando una función tiene más de 6 argumentos, los argumentos restantes se pasan en la pila tal que el argumento 7 se coloca abajo, el 8 encima, y así sucesivamente.
\subsection{Gestionando datos locales}
Hasta ahora, la mayoría de ejemplos que hemos visto no han requerido de variables locales, ya sea por optimización del compilador o por la sencillez del código,
sin embargo, sabemos que muchas veces es necesario almacenar variables locales.
Normalmente, en lenguajes basados en pila, o que soportan la recursión (e.g C, C++, Java, Rust, etc), las variables locales se almacenan en la pila.
Esta pila, se reserva en marcos (allocated in frames), que son bloques de memoria que contienen las variables locales de un procedimiento, estos marcos de memoria
se organizan tal que contienen información de retorno, almacenamiento local si es necesario y espacio temporal si es necesario, además hay un puntero de marco opcional:
$\%rbp$. Veamos un ejemplo de cadena de llamadas a procedimientos:
% \begin{center}
% 	\begin{tikzpicture}[>=Stealth, scale=1.5]
%
% 		\matrix (m1) [matrix of nodes, nodes={draw, minimum size=1cm}, column sep=0cm] {
% 			\node[fill=green!20] {\phantom{a}}; \\
% 			\node[fill=green!20] {\phantom{a}}; \\
% 			\node[fill=green!20] {\phantom{a}}; \\
% 			\node[fill=green!20] {\phantom{a}}; \\
% 		};
%
% 		\matrix (m2) [matrix of nodes, right=of m1, nodes={draw, minimum size=1cm}, column sep=0cm] {
% 			\node[fill=green!20] {\phantom{a}}; \\
% 			\node[fill=green!40] {fn A};                    \\
% 			\node[fill=green!20] {\phantom{a}}; \\
% 			\node[fill=green!20] {\phantom{a}}; \\
% 		};
%
% 		\draw[->] (m1) -- (m2);
%
% 		\matrix (m3) [matrix of nodes, right=of m2, nodes={draw, minimum size=1cm}, column sep=0cm] {
% 			\node[fill=green!20] {\phantom{a}}; \\
% 			\node[fill=green!40] {fn A};                    \\
% 			\node[fill=green!70] {fn B};                    \\
% 			\node[fill=green!0] {\phantom{a}};  \\
% 		};
%
% 		\draw[->] (m2) -- (m3);
% 		\matrix (m4) [matrix of nodes, right=of m3, nodes={draw, minimum size=1cm}, column sep=0cm] {
% 			\node[fill=green!20] {\phantom{a}}; \\
% 			\node[fill=green!40] {fn A};                    \\
% 			\node[fill=green!70] {fn B};                    \\
% 			\node[fill=green!70] {fn B};                    \\
% 		};
%
% 		\draw[->] (m3) -- (m4);
%
% 		\matrix (m5) [matrix of nodes, right=of m4, nodes={draw, minimum size=1cm}, column sep=0cm] {
% 			\node[fill=green!20] {\phantom{a}}; \\
% 			\node[fill=green!40] {fn A};                    \\
% 			\node[fill=green!70] {fn B};                    \\
% 			\node[fill=green!20] {\phantom{a}}; \\
% 		};
%
% 		\draw[->] (m4) -- (m5);
% 		\matrix (m6) [matrix of nodes, right=of m5, nodes={draw, minimum size=1cm}, column sep=0cm] {
% 			\node[fill=green!20] {\phantom{a}}; \\
% 			\node[fill=green!40] {fn A};                    \\
% 			\node[fill=green!20] {\phantom{a}}; \\
% 			\node[fill=green!20] {\phantom{a}}; \\
% 		};
%
% 		% Flecha entre la cuarta y quinta tabla
% 		\draw[->] (m5) -- (m6);
% 	\end{tikzpicture}
% \end{center}
Donde en cada uno de los pasos, $\%rsp$ apunta en la función más abajo existente y $\%rbp$ apunta al marco de la función inmediata anterior.
Los registros de programa actuan como un único recurso compartido por todos los procedimientos, aunque un solo procedimiento puede estar activo en un momento dado,
tenemos que asegurarnos de que cuando un procedimiento (caller) llama otro procedimiento (callee), el callee no sobreescriba los registros del caller. Veamos un ejemplo de esto:
\begin{center}
	\begin{minipage}{0.45\textwidth}
		\begin{minted}[fontsize=\large]{asm}
funcA:
    movq $15213, %rdx
    call funcB
    addq %rdx, %rax 
    ret
\end{minted}
	\end{minipage}
	\begin{minipage}{0.45\textwidth}
		\begin{minted}[fontsize=\large]{asm}
funcB: 
    subq $18213, %rdx
    ret 
    \end{minted}

	\end{minipage}
\end{center}
Como podemos ver, pudiera pasar que $\%rdx$ fuese sobreescrito en la función $funcB$, para evitar esto,
existen convenciones para coordinar el uso de registros.
Los registros $\%rbx, \%rbp, \%r12-\%r15$ se denominan \textit{callee-saved} ya que el callee debe guardar su valor si los modifica.
Para este fin, el callee puede o bien no modificar los registros, o guardando su valor en la pila y restaurándolo antes de salir, a este espacio se le llama
\textit{Saved Registers}, podemos verlo en el diagrama \ref{fig:stack}.
El resto de registros, excepto $\%rsp$, se clasifican como \textit{Caller-saved}, lo que
significa que pueden ser modificados por cualquier función, el nombre viene de un procedimiento $P$ que llama a un $Q$, como
$Q$ puede modificar libremente estos registros, $P$ debe guardar su valor si los necesita.
\subsection{Recursividad}
Las convenciones que hemos visto hasta ahora son suficientes para permitir a los procedimientos de x86-64 llamarse a sí mismos.
Cada llamada a procedimiento tiene su propio espacio en el stack, luego las variables locales de las múltiples llamadas no interfieren entre ellas.
Veamos un ejemplo del código C y ensamblador generado para una función recursiva:
\begin{center}
	\begin{minipage}{0.45\textwidth}
		\begin{minted}[fontsize=\large]{c}
long pcount_r(unsigned long x) {
    if (x == 0) {
        return 0;
    } else {
        return (x & 1) 
            + pcount_r(x >> 1);
    }
}
\end{minted}
	\end{minipage}
	\begin{minipage}{0.45\textwidth}
		\begin{minted}[fontsize=\large]{asm}
pcount_r:
    movl $0, %eax
    testq %rdi, %rdi
    je .L2
    pushq %rbx 
    movq %rdi, %rbx
    andl $1, %ebx
    shrq %rdi
    call pcount_r
    addq %rbx, %rax
    popq %rbx
  .L2
    rep; ret # Para branch prediction 
  \end{minted}
	\end{minipage}
\end{center}
Veamos más detenidamente lo que está pasando:
%Error 3
% \begin{itemize}
% 	\item Las primeras tres instrucciones sirven para saber cuando acaba la recursión, si x es 0, se devuelve 0.
% 	\item El pushq está ya que se va a modificar $\%rbx$ y se necesita restaurar su valor al salir de la función.
% 	\item Desde movq hasta shrq se calcula x $\&$ 1 y x >> 1.
% 	\item x >> 1 se pasa como argumento a la función $pcount_r$ (llamada recursiva)
% 	\item Se suma el resultado de la llamada recursiva con el valor de x \& 1.
% \end{itemize}
\section{Datos}
% Marker [DE AQUI HACIA ABAJO TODO CORRECTO]
\subsection{Arrays}
Los arrays en C se usan para almacenar colecciones de datos escalares en variables más grandes, la implementación particular
de C de los arrays es muy directa, una característica importante es que en C podemos manipular punteros a elementos de un array
y realizar operaciones con ellos. Los compiladores son extremadamente buenos simplificando el cálculo de direcciones 
usado cuando accedemos a un array, es por ello que la correspondencia entre C y ensamblador puede llegar a ser dificil de descifrar.
\subsubsection{Principios Básicos}
Sea $T$ un tipo de dato cualquiera, y $N$ una constante, consideramos la declaración $T A[N]$, entonces $A$ es un array de $N$ elementos.
La declaración reserva una región contigua de memoria de $T \cdot N$ bytes e introduce un identificador $A$ que se usa como puntero
al comienzo de esta región de memoria. Los elementos de este array se acceden mediante el uso de un índice entre 0 y N-1.
\begin{ejemplo}
    Consideramos \texttt{int arr[5]}, veamos los siguientes ejemplos:
    \begin{center}
        \includesvg[width=0.5\textwidth]{Imagenes/array}
        \end{center}
    \begin{itemize}
        \item \texttt{arr[4]} (int) se refiere al quinto elemento del array (valor 3).
        \item \texttt{arr} (int *) comienzo del array (posición arr)
        \item \texttt{arr+1} (int *) comienzo del array + 1 (posición arr + 4) 
        \item \texttt{\&arr[2]} (int *) dirección de val[2] (posición arr + 8)
        \item \texttt{arr[5]} (int) fuera de los límites del array, \href{https://en.cppreference.com/w/c/language/behavior}{undefined behavior.}
        \item \texttt{*(arr+1)} (int) valor de la dirección val + 1 (valor 5)
    \end{itemize}
\end{ejemplo}
Destacamos que declaramos arrays inmediatamente después de otro, no está garantizado que estos se almacenen en memoria de forma contigua. Así como hemos podido ver en el ejemplo anterior,
que C permite operaciones con los punteros, así, si $p$ es un puntero a un array de tipo $T$ y la dirección de memoria de $p$ es $x_p$, entonces la expresión $p + i$ tiene el valor \texttt{$x_p$ +  $T \cdot i$ = $i \cdot$ sizeof(T)}.
Recordamos que los operadores $\& $ y $*$ permiten la generación y dereferenciación de punteros respectivamente.
\begin{ejemplo}
    Ejemplo de acceso a elementos de array:
    \begin{center}
        \begin{minipage}{0.45\textwidth}
            \begin{minted}[fontsize=\large]{c}
int get_num(int *arr, int i) 
{
    return arr[i];
}
        \end{minted}
        \end{minipage}
        \begin{minipage}{0.45\textwidth}
        \begin{minted}[fontsize=\large]{asm}
    get_num:
        movslq  %esi, %rsi
        movl (%rdi, %rsi, 4), %eax
        ret
            \end{minted}
        \end{minipage}
    \end{center}
    Recordamos que \texttt{movslq} es una instrucción que significa \textit{Move with Sign-extend Long to Quad}, se utiliza ya que para hacer 
    operaciones con la dirección de memoria, es necesario que los registros sean de 64 bits.
\end{ejemplo}

\subsubsection{Arrays Anidados}
Los principios de los arrays multidimensionales son los mismos que los de los arrays unidimensionales, de esta manera, la declaración
\texttt{T A[R][C]} declara un array de $R$ elementos, cada uno de los cuales es un array de $C$ elementos. En general, un array bidimensional de elementos de tipo T. De otra forma,
podemos decir que tiene $R$ filas y $C$ columnas. Notemos que si el tamaño de $T$ es $K$ bytes, entonces el tamaño de $A$ es $ C \cdot R \cdot K$ bytes.
La disposición de los elementos en memoria es por filas, \textit{(row-major-order)}:
\begin{center}
\includesvg[width=0.8\textwidth]{Imagenes/array2d.svg}
\end{center}
De esta forma, si tenemos un array \texttt{T A[R][C]} el elemento 
\texttt{T[i][j]} está en la posición de memoria \texttt{$\&$A[i][j] = $x_A$ + L(C $\cdot i + j$)}.
\begin{ejemplo}
    Veamos un ejemplo en el que el flag --no-pie debe ser considerado
    \begin{center}
        \begin{minipage}{0.45\textwidth}
            \begin{minted}[fontsize=\large]{c}
    int *get_row(int index) 
    {
        return arr[index];
    }
            \end{minted}
        \end{minipage}
        \begin{minipage}{0.45\textwidth}
            \begin{minted}[fontsize=\large]{c}
#define ROWS 4
// Zip_dig -> Array de 5 elementos
zip_dig arr[ROWS] = {
    {1, 2, 7, 4, 5},
    {6, 7, 2, 9, 1},
    {1, 9, 3, 1, 5},
    {1, 7, 8, 9, 2}};
            \end{minted}
        \end{minipage}
        \vspace{0.5cm}
        \vspace{0.5cm}
        \begin{minipage}{0.45\textwidth}
            \begin{minted}[fontsize=\large]{asm}
    get_num:
        leaq (%rdi, %rdi, 4), %rdx
        leaq arr(,%rdx, 4), %rax
        ret
            \end{minted}
        \end{minipage}
        \begin{minipage}{0.45\textwidth}
            \begin{minted}[fontsize=\large]{asm}
    get_num:
        movslq %edi, %rdi
        leaq arr(%rdi ,%rdi, 4), %rdx
        leaq arr(%rip), %rax
        leaq (%rax, %rdx, 4), %rax
        ret
        \end{minted}
        \end{minipage}
    \end{center}
    A la izquierda, como se usa --no-pie que desactiva el \textit{Position Independent Code}, la dirección de memoria es conocida en tiempo de compilación, por lo que
    se puede acceder directamente con $ 5 \cdot \text{index}$ por ser 5 elementos por fila y $4 \cdot \text{index}$ por ser un tipo de dato entero de 4 bytes.
    En el segundo caso, el direccionamiento es relativo a contador de programa $\%rip$. Notemos que --no-pie resulta más eficiente sin embargo es más vulnerable a ataques de buffer overflow.
\end{ejemplo}
\subsubsection{Array Multinivel}
De manera similar a los arrays bidimensionales, tenemos los arrays multinivel, que tienen la peculiaridad de que 
se tiene un array de punteros a arrays, de esta forma, para acceder a un dato son necesarios dos accesos a memoria, uno para acceder al puntero y otro para acceder al dato.
\begin{ejemplo}
    Un ejemplo de array multinivel en C: 
    \begin{center}
        \begin{minted}{c}
zip_dig cmu = {0, 5, 2, 7, 3};
zip_dig mit = {1, 2, 3, 4, 5};
zip_dig ucb = {9, 3, 7, 5, 1};

# define UCOUNT 3
int *univ[UCOUNT] = {cmu, mit, ucb};
        \end{minted}
    \end{center}
\end{ejemplo}
\subsection{Estructuras}
C permite dos formas de definir tipos de datos compuestos: \textit{struct} que agrupan múltiples datos en un 
solo objeto y \textit{union} que permiten que un objeto sea referenciado usando diferentes tipos de datos. Veamos en detalle las estructuras.
\\ \\ 
El struct en C crea un tipo de dato que agrupa múltiples variables que pueden ser de diferentes tipos en un solo objeto,
los distintos elementos de la estructura se referencian mediante nombres, con la sintaxis \texttt{struct.field, struct->field}.
El struct se parece a los arrays en el sentido de que se garantiza que los datos de un mismo struct están contiguos en memoria.
Es el compilador el que mantiene información sobre cada estructura para indicar el offset de cada campo. El tamaño de los 
structs es el tamaño de todos sus campos sumados, teniendo en cuenta que es posible que el compilador reserve espacio adicional tras un 
tipo de dato para alinear la memoria.
\begin{ejemplo}
    Veamos un ejemplo de struct en C y cómo se distribuye su memoria:
    \begin{center}
        \begin{minipage}{0.45\textwidth}
            \begin{minted}[fontsize=\large]{c}
        struct rec {
            int a[4];
            size_t i;
            struct rec *next;
        }
            \end{minted}
        \end{minipage}
        \begin{minipage}{0.48\textwidth}
            \includesvg[width=0.8\textwidth]{Imagenes/struct}
        \end{minipage}
    \end{center}
        Notemos que si en el struct el array de enteros fuese de tamaño 3, el tamaño del struct se mantendría en 32
        B por el alineamiento de la memoria. El puntero $r$ indica el primer byte de la estructura, se accede
        mediante el offset, que como hemos dicho lo determina el compilador en tiempo de compilación.
        \begin{center}
        \begin{minipage}{0.48\textwidth}
            \begin{minted}[fontsize=\large]{c}
    void set_i(struct rec *r, size_t val)
    {
        r->i = val;
    }
            \end{minted}
        \end{minipage}
        \begin{minipage}{0.45\textwidth}
            \begin{minted}[fontsize=\large]{asm}
        set_i:
            movq %rsi, (%rdi, 16)
            ret
            \end{minted} 
        \end{minipage}
    \end{center}
\end{ejemplo}
Hemos introducido ya el alineamiento de los datos, detallamos los principios que rigen su funcionamiento:
\begin{itemize}
    \item Si el tipo de dato primitivo requiere $K$ bytes, la dirección debe ser múltiplo de $K$.
    \item El alineamiento es requisito en algunas maquinas, en x86-64 es recomendado por motivos de eficiencia (Es ineficiente cargar o almacenar datos que cruzan fronteras de quad word y
        la memoria virtual es delicada cuando un dato se extiende a 2 páginas).
    \item En la mayoría de lenguajes compilados (C, C++, Rust, etc) es el compilador el que inserta huecos en la estructura para alinear los campos.
    \item El alineamiento se hace:
        \begin{itemize}
            \item Para el struct, se toma el máximo valor miembro para alinear, esto se llama alineamiento externo.
            \item Para los datos primitivos, se toma el tamaño del dato, es decir se añade padding antes del dato miembro si fuese necesario para que se encuentre en una posición múltiplo de su tamaño.
        \end{itemize}
\end{itemize}
El compilador no reordena los datos con el fin de ahorrar espacio, sin embargo
el programador debería considerar siempre cómo ordenar los datos para ocupar el menor espacio posible, esto es 
especialmente importante ya que a diferencia de otras optimizaciones que el compilador hace por nosotros, esta no se puede hacer ya que rompería la coherencia y predictibilidad del código.
\subsection{Uniones}
Las uniones proporcionan una manera de circunvalar el sistema de tipos de C, permitiendo a un solo objeto ser referenciado por múltiples tipos. La sintaxis de la unión es idéntica a la de un struct, pero 
sus semánticas son diferentes. En una unión, todos los campos comparten la misma dirección de memoria, lo que significa que si se modifica un campo, se modifica el resto, para hacer esto, el compilador reserva
tamaño de acuerdo al elemento más grande.
\begin{ejemplo}
    Veamos un ejemplo de una unión en comparación con un struct:
    \begin{center}
        \begin{minipage}{0.45\textwidth}
            \begin{minted}[fontsize=\large]{c}
    union U1 {
        char c;
        int i[2];
        double v;
    } *up;
    \end{minted}
    \end{minipage}
    \begin{minipage}{0.45\textwidth}
        \begin{minted}[fontsize=\large]{c}
    struct S1 {
        char c;
        int i[2];
        double v;
    } *sp;
        \end{minted}
    \end{minipage}
\end{center} 
Si compilamos ambos códigos en una máquina Linux en x86-64, el offset y tamaño total de las estructuras sería el siguiente:
\begin{center}
    \begin{tabular}{ccccc}
        Tipo & c & i & v & Total Size \\ 
        \hline 
        U1 & 0 & 0 & 0 & 8 \\ 
        S1 & 0 & 4 & 16 & 24 \\ 
    \end{tabular}
\end{center}
Este gran poder conlleva una gran responsabilidad, ya que el programador debe ser consciente de que 
si usa datos como \texttt{float} e \texttt{int} en la misma unión, 
e intenta interpretar los bits del entero como si fuese un float,
el resultado no será el esperado, ya que no tendremos el mismo efecto que el cast (float) sobre el entero.
\end{ejemplo}
Recordamos que las palabras se almacenan en memoria como un número N de bytes, en esta organización, el byte mas 
significativo puede ser el que está en la dirección más baja (Big Endian) o en la dirección más alta (Little Endian).
Esta diferencia puede causar problemas al intercambiar datos binarios entre máquinas, también mencionamos Middle endian, que son  
aquellas arquitecturas que se pueden configurar de ambas maneras, aparece en ARM, sin embargo, la mayoría de las máquinas modernas son Little Endian.
\\ \\ 
Para finalizar este capítulo, presentamos un resumen de los tipos compuestos en C:
\begin{itemize}
    \item \textbf{Arrays}
        \begin{itemize}
            \item Reservan memoria contigua para almacenar elementos
            \item Se usa aritmética de indexación para localizar elementos individuales
            \item Puntero al primer elemento
            \item No comprueban límites
        \end{itemize}
\item \textbf{Structs}
    \begin{itemize}
        \item Reserva de una sola región de memoria, los campos en el mismo orden declarado 
        \item Se accede usando desplazamientos determinados por el compilador
        \item Alineamiento interno y externo mediante relleno
    \end{itemize}
\item \textbf{Unions}
    \begin{itemize}
        \item Declaraciones superpuestas
        \item Proveen de una manera para saltarse el sistema de tipos de C
    \end{itemize}
\end{itemize}


    \fancyhead[R]{\helv \nouppercase{\rightmark}}
    \chapter{Relaciones de Problemas}
    \section{Introducción}

\begin{ejercicio}
    Considerar el siguiente fragmento de programa para 2 procesos \verb|P1| y \verb|P2|: Los dos procesos
    pueden ejecutarse a cualquier velocidad. ¿Cuáles son los posibles valores resultantes para la
    variable \verb|x|? Suponer que \verb|x| debe ser cargada en un registro para incrementarse y que cada
    proceso usa un registro diferente para realizar el incremento.
    \setlength{\columnsep}{2cm} % Ajusta el espacio entre columnas
    \begin{multicols}{2}
        \begin{minted}{pascal}
        {variables compartidas}
        var x : integer := 0 ;
        Process P1;
        var i: integer;
        begin
          begin
            for i:= 1 to 2 do begin
              x:= x + 1;
            end
          end
        end
        \end{minted}
        
        \begin{minted}{pascal}
            

        Process P2;
        var j: integer;
        begin
          begin
            for j:= 1 to 2 do begin
              x:= x + 1;
            end
          end
        end
        \end{minted}
    \end{multicols}

Observando el código, cada proceso hace 2 lecturas y dos escrituras (incrementos) en \verb|x|.
\begin{itemize}
    \item Como cada proceso aumenta dos veces el valor de \verb|x|, el valor de \verb|x| ha de ser, como mínimo, $2$.
    \item Como en total se hacen 4 incrementos, el valor de \verb|x| ha de ser $4$ como máximo.
\end{itemize}

Notando por $l_{ij}$ a la $j$-ésima lectura del proceso $i$ y por $e_{ij}$ a la $j$-ésima escritura del proceso $i$, ambas referidas a la variable \verb|x|, podemos obtener cualquiera de las siguientes trazas de ejecución:

\begin{table}[H]
\centering
\begin{tabular}{|c|c|c||c|c|c||c|c|c||c|c|c|}
    \hline
    \verb|P1| & \verb|P2| & \verb|x| & \verb|P1| & \verb|P2| & \verb|x| & \verb|P1| & \verb|P2| & \verb|x| & \verb|P1| & \verb|P2| & \verb|x| \\
    \hline
    \hline
    $l_{11}$ & - & 0 & $l_{11}$ & - & 0 & $l_{11}$ & - & 0 & $l_{11}$ & - & 0 \\
    \hline
    $e_{11}$ & - & 1 & - & $l_{21}$ & 0 & $e_{11}$ & - & 1 & - & $l_{21}$ & 0 \\
    \hline
    - & $l_{21}$ & 1 & $e_{11}$ & - & 1 & - & $l_{21}$ & 1 & $e_{11}$ & - & 1 \\
    \hline
    - & $e_{21}$ & 2 & - & $e_{21}$ & 1 & - & $e_{21}$ & 2 & - & $e_{21}$ & 1 \\
    \hline
    $l_{12}$ & - & 2 & $l_{12}$ & - & 1 & $l_{12}$ & - & 2 & $l_{12}$ & - & 1 \\
    \hline
    $e_{12}$ & - & 3 & $e_{12}$ & - & 2 & - & $l_{22}$ & 2 & - & $l_{22}$ & 1 \\
    \hline
    - & $l_{22}$ & 3 & - & $l_{22}$ & 2 & $e_{12}$ & - & 3 & $e_{12}$ & - & 2 \\
    \hline
    - & $e_{22}$ & 4 & - & $e_{22}$ & 3 & - & $e_{22}$ & 3 & - & $e_{22}$ & 2  \\
    \hline
\end{tabular}
\end{table}

Luego los posibles valores resultantes para \verb|x| son: 2, 3 y 4.

\end{ejercicio}


\begin{ejercicio}
    ¿Cómo se podría hacer la copia del fichero \verb|f| en otro \verb|g|, de forma concurrente, utilizando la
    instrucción concurrente \verb|cobegin-coend|? Para ello, suponer que:
    \begin{enumerate}
        \item Los archivos son una secuencia de items de un tipo arbitrario \verb|T|, y se encuentran ya abiertos
        para lectura (\verb|f|) y escritura (\verb|g|). Para leer un ítem de \verb|f| se usa la llamada a función \verb|leer(f)| y
        para saber si se han leído todos los ítems de \verb|f|, se puede usar la llamada \verb|fin(f)| que devuelve
        verdadero si ha habido al menos un intento de leer cuando ya no quedan datos. Para
        escribir un dato \verb|x| en \verb|g| se puede usar la llamada a procedimiento \verb|escribir(g,x)|.

        \item El orden de los items escritos en \verb|g| debe coincidir con el de \verb|f|.
        \item Dos accesos a dos archivos distintos pueden solaparse en el tiempo.
    \end{enumerate}~\\

\noindent
La copia del fichero \verb|f| en el fichero \verb|g| se podría realizar siguiendo el paradigma productor/consumidor que hemos visto en teoría en el Tema 1, mediante el uso de dos procesos:
\begin{itemize}
    \item Uno que lea un ítem del fichero \verb|f| y lo escriba en una variable compartida.
    \item Otro que lea dicha variable compartida y escriba el ítem en el fichero \verb|g|.
\end{itemize}
En dicho código, debemos garantizar que:
\begin{itemize}
    \item El consumidor no lea la variable antes de que el productor escriba en ella.
    \item En la segunda escritura del productor, debemos esperar a que antes la haya leído el consumidor.
    \item En la segunda lectura del consumidor, debemos esperar a que antes haya modificado la variable el productor.
\end{itemize}
Siguiendo estos pasos, obtendríamos un código como el siguiente:
\begin{minted}{pascal}
    process CopiaFicheros ;
    var ant, sig : T ;
    begin
      sig = leer(f) ;
      while not fin(f) do begin
        ant = sig ;
        cobegin
          escribir(g, anterior) ;
          sig = leer(f) ;
        coend
      end
    end
\end{minted}


\end{ejercicio}

\begin{ejercicio}\label{ej:3}
    Construir, utilizando las instrucciones concurrentes \verb|cobegin-coend| y \verb|fork-join|, programas concurrentes que se correspondan con los grafos de precedencia que se muestran en la figura~\ref{fig:grafoEj3}.
    \begin{figure}
        \centering
        \begin{subfigure}{0.3\textwidth}
            \centering
            \resizebox{\linewidth}{!}{
                \begin{tikzpicture}[
                    node/.style={circle, draw, minimum size=0.5cm},
                    edge/.style={-stealth}
                    ]
        
                    % Nodos
                    \node[node] (P0) {P0};
                    \node[node, below left=of P0] (P1) {P1};
                    \node[node, below right=of P0] (P2) {P2};
                    \node[node, below=of P1] (P3) {P3};
                    \node[node, below left=of P3] (P4) {P4};
                    \node[node, below right=of P3] (P5) {P5};
                    \node[node, below=of P5] (P6) {P6};
        
                    % Aristas
                    \draw[edge] (P0) -- (P1);
                    \draw[edge] (P0) -- (P2);
                    \draw[edge] (P1) -- (P3);
                    \draw[edge] (P3) -- (P4);
                    \draw[edge] (P3) -- (P5);
                    \draw[edge, bend right] (P4) to (P6);
                    \draw[edge] (P5) -- (P6);
                    \draw[edge, bend left] (P2) to (P6);
                
                \end{tikzpicture}
            }
            \caption{DAG del apartado~\ref{ej:3.1}.}
            \label{fig:grafoEj3.1}
            
        \end{subfigure}
        \begin{subfigure}{0.3\textwidth}
            \centering
            \resizebox{\linewidth}{!}{
                \begin{tikzpicture}[
                    node/.style={circle, draw, minimum size=1cm},
                    edge/.style={-stealth}
                    ]
        
                    % Nodos
                    \node[node] (P0) {P0};
                    \node[node, below left=of P0] (P1) {P1};
                    \node[node, below right=of P0] (P2) {P2};
                    \node[node, below=of P1] (P3) {P3};
                    \node[node, below left=of P1] (P4) {P4};
                    \node[node, below right=of P1] (P5) {P5};
                    \node[node, below=of P5] (P6) {P6};
        
                    % Aristas
                    \draw[edge] (P0) -- (P1);
                    \draw[edge] (P0) -- (P2);
                    \draw[edge] (P1) -- (P3);
                    \draw[edge] (P1) -- (P4);
                    \draw[edge] (P1) -- (P5);
                    \draw[edge, bend right] (P4) to (P6);
                    \draw[edge] (P5) -- (P6);
                    \draw[edge, bend left] (P2) to (P6);
                    \draw[edge, bend right] (P3) to (P6);
                
                \end{tikzpicture}
            }
            \caption{DAG del apartado~\ref{ej:3.2}.}
            \label{fig:grafoEj3.2}
            
        \end{subfigure}
        \begin{subfigure}{0.3\textwidth}
            \centering
            \resizebox{\linewidth}{!}{
                \begin{tikzpicture}[
                    node/.style={circle, draw, minimum size=1cm},
                    edge/.style={-stealth}
                    ]
        
                    % Nodos
                    \node[node] (P0) {P0};
                    \node[node, below left=of P0] (P1) {P1};
                    \node[node, below right=of P0] (P2) {P2};
                    \node[node, below=of P1] (P3) {P3};
                    \node[node, below left=of P3] (P4) {P4};
                    \node[node, below right=of P3] (P5) {P5};
                    \node[node, below=of P5] (P6) {P6};
        
                    % Aristas
                    \draw[edge] (P0) -- (P1);
                    \draw[edge] (P0) -- (P2);
                    \draw[edge] (P1) -- (P3);
                    \draw[edge] (P3) -- (P4);
                    \draw[edge] (P3) -- (P5);
                    \draw[edge, bend right] (P4) to (P6);
                    \draw[edge] (P5) -- (P6);
                    \draw[edge] (P2) to (P5);
                
                \end{tikzpicture}
            }
            \caption{DAG del apartado~\ref{ej:3.3}.}
            \label{fig:grafoEj3.3}
            
        \end{subfigure}
        \caption{Grafos de precedencia del ejercicio~\ref{ej:3}.}
        \label{fig:grafoEj3}
    \end{figure}


    \begin{enumerate}
        \item \label{ej:3.1}
         Grafo de precedencia de la figura~\ref{fig:grafoEj3.1}:

         \begin{figure}[H]
             \centering
             \begin{subfigure}[b]{0.45\textwidth}
                \centering
                 \begin{minted}{pascal}
                     begin
                       P0;
                       fork P2; P1;
                       P3;
                       fork P5; P4;
                       join P2; join P5;
                       P6;
                     end
                 \end{minted}
             \end{subfigure}\hfill
             \begin{subfigure}[b]{0.45\textwidth}
                \centering
                 \begin{minted}{pascal}
                     begin
                       P0;
                       cobegin
                         P2;
                         begin
                           P1;
                           P3;
                           cobegin P4; P5; coend
                         end
                       coend
                       P6;
                     end
                 \end{minted}
             \end{subfigure}
         \end{figure}
                 
        \item \label{ej:3.2}
        Grafo de precedencia de la figura~\ref{fig:grafoEj3.2}:
        \begin{figure}[H]
            \centering
            \begin{subfigure}[b]{0.45\textwidth}
                \centering
                \begin{minted}{pascal}
                begin
                  P0;
                  fork P2; P1;
                  fork P5; fork P3; P4;
                  join P2; join P5; join P3;
                  P6;
                end
                \end{minted}
            \end{subfigure}\hfill
            \begin{subfigure}[b]{0.45\textwidth}
                \centering
                \begin{minted}{pascal}
                begin
                  P0;
                  cobegin
                    P2;
                    begin
                      P1;
                      cobegin P4; P3; P5; coend
                    end
                  coend
                  P6;
                end
                \end{minted}
            \end{subfigure}
        \end{figure}
        
         \item \label{ej:3.3}
         Grafo de precedencia de la figura~\ref{fig:grafoEj3.3}:
         \begin{minted}{pascal}
         begin
           P0;
           fork P2; P1;
           P3;
           fork P4; join P2; P5;
           join P4;
           P6;
         end
         \end{minted}

         Sin embargo, no podemos hacer al 100\% el DAG de la figura~\ref{fig:grafoEj3.3}, ya que tras \verb|P3| debemos crear una estructura \verb|cobegin-coend|. Sin embargo, este debe esperar a \verb|P2|, por lo que la estructura \verb|cobegin-coend| tendrá que esperar a \verb|P2|, pero es que \verb|P4| no necesita que \verb|P2| termine. 

         Por tanto, no se puede programar con creación de hebras de forma estructurada. Sin embargo, podemos ofrecer dos soluciones, cada una que impone algo que el grafo no nos dice:
         \begin{enumerate}
             \item Si obligamos a que \verb|P4| también espere a \verb|P2|, obtendríamos el código:
                 \begin{minted}{pascal}
                     begin
                       P0;
                       cobegin
                         P2;
                         begin
                           P1; P3;
                         end
                       coend
                       cobegin P4; P5; coend
                       P6;
                     end
                 \end{minted}
                 
             \item Si ahora queremos ejecutar de forma concurrente el flujo que tiene a \verb|P1|, \verb|P3| y \verb|P4| con el flujo que tiene a \verb|P2|, entonces obligamos a que \verb|P5| espere a \verb|P4| (que no nos lo especifica el DAG, pero lo necesitamos para poder programarlo de forma estructurada):
                 \begin{minted}{pascal}
                     begin
                       P0;
                       cobegin
                         begin P1; P3; P4; end
                         P2;
                       coend
                       P5;
                       P6;
                     end
                 \end{minted}
                 
         \end{enumerate}
    \end{enumerate}

\end{ejercicio}



\begin{ejercicio} \label{ej:4}
    Dados los siguientes fragmentos de programas concurrentes, obtener sus grafos de precedencia asociados:
    \begin{figure}[H]
        \centering
        \begin{subfigure}[b]{0.45\textwidth}
            \centering
            \begin{minted}{pascal}
                begin
                    P0 ;
                    cobegin
                        P1 ;
                        P2 ;
                        cobegin
                            P3 ; P4 ; P5 ; P6 ;
                        coend ;
                        P7 ;
                    coend
                    P8 ;
                end
            \end{minted}
            \caption{Programa 1.}
            \label{code:prog1_Ej4}
        \end{subfigure}\hfill
        \begin{subfigure}[b]{0.45\textwidth}
            \centering
            \begin{minted}{pascal}
                begin
                    P0 ;
                    cobegin
                        begin
                            cobegin
                                P1 ; P2 ;
                            coend
                            P5 ;
                        end
                        begin
                            cobegin
                                P3 ; P4 ;
                            coend
                            P6 ;
                        end
                    coend
                    P7 ;
                end
            \end{minted}
            \caption{Programa 2.}
            \label{code:prog2_Ej4}
        \end{subfigure}
        \caption{Programas concurrentes del ejercicio~\ref{ej:4}.}
    \end{figure}
    
    \begin{enumerate}
        \item Programa de la figura~\ref{code:prog1_Ej4}.

            \begin{figure}[H]
            \centering
            \begin{tikzpicture}[
                    node/.style={circle, draw, minimum size=1cm},
                    edge/.style={-stealth}
                    ]
                \node[node] (P0) {P0};
                \node[node, below=of P0] (P4) {P4};
                \node[node, left=of P4] (P3) {P3};
                \node[node, left=of P3] (P2) {P2};
                \node[node, left=of P2] (P1) {P1};
                \node[node, right=of P4] (P5) {P5};
                \node[node, right=of P5] (P6) {P6};
                \node[node, right=of P6] (P7) {P7};
                \node[node, below=of P4] (P8) {P8};

                \draw[edge] (P0) -- (P1);
                \draw[edge] (P0) -- (P2);
                \draw[edge] (P0) -- (P3);
                \draw[edge] (P0) -- (P4);
                \draw[edge] (P0) -- (P5);
                \draw[edge] (P0) -- (P6);
                \draw[edge] (P0) -- (P7);

                \draw[edge] (P1) -- (P8);
                \draw[edge] (P2) -- (P8);
                \draw[edge] (P3) -- (P8);
                \draw[edge] (P4) -- (P8);
                \draw[edge] (P5) -- (P8);
                \draw[edge] (P6) -- (P8);
                \draw[edge] (P7) -- (P8);
            \end{tikzpicture}
            \caption{DAG para la figura~\ref{code:prog1_Ej4}.}
            \end{figure}
        \item Programa de la figura~\ref{code:prog2_Ej4}.
            \begin{figure}[H]
            \centering
            \begin{tikzpicture}[node/.style={circle, draw, minimum size=1cm},
                                edge/.style={-stealth}]
                \node[node] (P0) {P0};
                \node[node, below left=of P0] (P2) {P2};
                \node[node, left=of P2] (P1) {P1};
                \node[node, below right=of P0] (P3) {P3};
                \node[node, right=of P3] (P4) {P4};
                \node[node, below right=of P1] (P5) {P5};
                \node[node, below left=of P4] (P6) {P6};

                \node[below=of P0] (A) {};
                \node[below=of A] (B) {};
                \node[node, below=of B] (P7) {P7};
                
                \draw[edge] (P0) -- (P1);
                \draw[edge] (P0) -- (P2);
                \draw[edge] (P0) -- (P3);
                \draw[edge] (P0) -- (P4);
                \draw[edge] (P1) -- (P5);
                \draw[edge] (P2) -- (P5);
                \draw[edge] (P3) -- (P6);
                \draw[edge] (P4) -- (P6);
                \draw[edge] (P5) -- (P7);
                \draw[edge] (P6) -- (P7);
            \end{tikzpicture}
            \caption{DAG para la figura~\ref{code:prog2_Ej4}.}
            \end{figure}
    \end{enumerate}
\end{ejercicio}

\begin{ejercicio} \label{ej:5}
    Suponer un sistema de tiempo real que dispone de un captador de impulsos conectado a un
    contador de energía eléctrica. La función del sistema consiste en contar el número de impulsos
    producidos en 1 hora (cada Kwh consumido se cuenta como un impulso) e imprimir este número
    en un dispositivo de salida. Para ello se dispone de un programa concurrente con 2 procesos: un
    proceso acumulador (lleva la cuenta de los impulsos recibidos) y un proceso escritor (escribe
    en la impresora). En la variable común a los 2 procesos \verb|n| se lleva la cuenta de los impulsos. El
    proceso acumulador puede invocar un procedimiento \verb|Espera_impulso| para esperar a que llegue
    un impulso, y el proceso escritor puede llamar a \verb|Espera_fin_hora| para esperar a que termine
    una hora. El código de los procesos de este programa podría ser el descrito en el Código Fuente~\ref{code:ej5}.
    \begin{observacion}
        En el programa se usan sentencias de acceso a la variable \verb|n| encerradas entre los símbolos \verb|<| y
        \verb|>|. Esto significa que cada una de esas sentencias se ejecuta en exclusión mutua entre los dos
        procesos, es decir, esas sentencias se ejecutan de principio a fin sin entremezclarse entre ellas.
        Supongamos que en un instante dado el acumulador está esperando un impulso, el escritor está
        esperando el fin de una hora, y la variable \verb|n| vale \verb|k|. Después se produce de forma simultánea
        un nuevo impulso y el fin del periodo de una hora.
    \end{observacion}

    Obtener las posibles secuencias de interfolicación de las instrucciones (1),(2), y (3) a partir de
    dicho instante, e indicar cuales de ellas son correctas y cuales incorrectas (las incorrectas son
    aquellas en las cuales el impulso no se contabiliza).
    \begin{listing}
        \begin{minted}{pascal}
            { variable compartida: }
            var n : integer; { contabiliza impulsos }
            begin
            while true do begin
                Espera_impulso();
                < n := n+1 > ; { (1) }
                end
            end
            process Escritor ;
            begin
            while true do begin
                Espera_fin_hora();
                write( n ) ; { (2) }
                < n := 0 > ; { (3) }
                end
            end
        \end{minted}
        \caption{Código acumulador-escritor del ejercicio~\ref{ej:5}.}
        \label{code:ej5}
    \end{listing}

    En primer lugar, notemos que la intrucción $2$ siempre se ejecutará
    antes que la instrucción $3$, ya que están ambas en el mismo proceso (el escritor).
    Por tanto, las secuencia de instrucciones que se pueden dar son las intercalaciones de la instrucción $1$ con las otras dos instrucciones; es decir: $A)~1-2-3$, $B)~2-1-3$ y $C)~2-3-1$. El análisis de cada una de ellas está en la Tabla~\ref{tab:ej5}.
    \begin{table}
        \centering
        \begin{tabular}{|lcc|lcc|lcc|}
            \hline
            \multicolumn{3}{|c|}{Opción $A$} & \multicolumn{3}{|c|}{Opción $B$} & \multicolumn{3}{|c|}{Opción $C$} \\ \hline
            Operación & \verb|n| & Salida & Operación & \verb|n| & Salida & Operación & \verb|n| & Salida\\
            \hline
            $-$ & $k$ & $-$ & $-$ & $k$ & $-$ & $-$ & $k$ & $-$ \\
            \verb|1|$\to$ \verb|n:=n+1| & $k+1$ & $-$ & \verb|2|$\to$ \verb|write(n)| & $k$ & $k$ & \verb|2|$\to$ \verb|write(n)| & $k$ & $k$ \\
            \verb|2|$\to$ \verb|write(n)| & $k+1$ & $k+1$ & \verb|1|$\to$ \verb|n:=n+1| & $k+1$ & $k$ & \verb|3|$\to$ \verb|n:=0| & $0$ & $k$ \\
            \verb|3|$\to$ \verb|n:=0| & $0$ & $k+1$ & \verb|3|$\to$ \verb|n:=0| & $0$ & $k$ & \verb|1|$\to$ \verb|n:=n+1| & $1$ & $k$ \\
            \hline
        \end{tabular}
        \caption{Tabla de opciones del ejercicio~\ref{ej:5}.}
        \label{tab:ej5}
    \end{table}
    Notemos que tanto la opción $A$ como la $C$ son válidas, ya que en ambas se contabiliza el impulso. No obstante, difieren entre sí, puesto que en la opción $A$ se contabiliza el impulso en la hora que termina (imprimiéndose entonces), mientras que en la opción $C$ se contabiliza el impulso en la hora siguiente (no imprimiéndose entonces en esta salida). No obstante, la opción $B$ es incorrecta, ya que no se contabiliza el impulso en la hora que termina ni en la siguiente.
\end{ejercicio}



\begin{ejercicio} \label{ej:6}
    Supongamos un programa concurrente en el cual hay, en memoria compartida dos vectores \verb|a| y
    \verb|b| de enteros y con tamaño par, declarados como sigue:
    \begin{minted}{pascal}
        var a,b : array[1..2*n] of integer ; { n es una constante predefinida }
    \end{minted}
    Queremos escribir un programa para obtener en \verb|b| una copia ordenada del contenido de \verb|a| (nos
    da igual el estado en que queda \verb|a| después de obtener \verb|b|). Para ello disponemos de la función
    \verb|Sort| que ordena un tramo de \verb|a| (entre las entradas \verb|s| y \verb|t|, ambas incluidas). También disponemos
    la función \verb|Copiar|, que copia un tramo de \verb|a| (desde \verb|s| hasta \verb|t|) en \verb|b| (a partir de \verb|o|). Estas funciones
    se muestran en el Código Fuente~\ref{code:ej_6SortCopiar}.
    \begin{listing}
        \begin{minted}{pascal}
            procedure Sort( s,t : integer );
                var i, j : integer ;
                begin
                    for i := s to t do
                    for j:= s+1 to t do
                        if a[i] < a[j] then
                            swap( a[i], b[j] ) ;
                end

            procedure Copiar( o,s,t : integer );
                var d : integer ;
                begin
                    for d := 0 to t-s do
                        b[o+d] := a[s+d] ;
                end
        \end{minted}
        \caption{Procedimientos \mintinline{pascal}{Sort} y \mintinline{pascal}{Copiar} del ejercicio~\ref{ej:6}.}
        \label{code:ej_6SortCopiar}
    \end{listing}

    El programa para ordenar se puede implementar de dos formas:
    \begin{enumerate}
        \item Ordenar todo el vector \verb|a|, de forma secuencial con la función \verb|Sort|, y después copiar cada
        entrada de \verb|a| en \verb|b|, con la función \verb|Copiar|.
        \item Ordenar las dos mitades de \verb|a| de forma concurrente, y después mezclar dichas dos mitades
        en un segundo vector \verb|b| (para mezclar usamos un procedimiento \verb|Merge|).
    \end{enumerate}
    En el Código Fuente~\ref{code:ej6_2versiones} se muestra el código de ambas versiones.
    \begin{listing}
        \begin{minted}{pascal}
            procedure Secuencial() ;
                var i : integer ;
                begin
                    Sort( 1, 2*n ); { ordena a }
                    Copiar( 1, 2*n ); { copia a en b }
                end

            procedure Concurrente() ;
                begin
                    cobegin
                        Sort( 1, n );
                        Sort( n+1, 2*n );
                    coend
                    Merge( 1, n+1, 2*n );
                end
        \end{minted}
        \caption{Procedimientos \mintinline{pascal}{Secuencial} y \mintinline{pascal}{Concurrente} del ejercicio~\ref{ej:6}.}
        \label{code:ej6_2versiones}
    \end{listing}

    El código de la función \verb|Merge|, disponible en el Código Fuente~\ref{code:ej6_Merge}, se encarga de ir leyendo las dos mitades de \verb|a|, en cada paso, seleccionar el menor elemento de los dos siguientes por leer (uno en cada mitad), y escribir dicho menor elemento en la siguiente mitad del vector mezclado \verb|b|.
    \begin{listing}
        \begin{minted}{pascal}
            procedure Merge( inferior, medio, superior: integer ) ;
                { siguiente posicion a escribir en b }
                var escribir : integer := 1 ;
                { siguiente pos. a leer en primera mitad de a }
                var leer1 : integer := inferior ;
                { siguiente pos. a leer en segunda mitad de a }
                var leer2 : integer := medio ;
                begin
                    { mientras no haya terminado con alguna mitad }
                    while leer1 < medio and leer2 <= superior do begin
                        if a[leer1] < a[leer2] then begin { minimo en la primera mitad }
                            b[escribir] := a[leer1] ;
                            leer1 := leer1 + 1 ;
                        end else begin { minimo en la segunda mitad }
                            b[escribir] := a[leer2] ;
                            leer2 := leer2 + 1 ;
                        end
                        escribir := escribir+1 ;
                    end
                    { se ha terminado de copiar una de las mitades,
                    copiar lo que quede de la otra }
                    if leer2 > superior then
                        { copiar primera } Copiar( escribir, leer1, medio-1 );
                    else Copiar( escribir, leer2, superior ); { copiar segunda }
                end
        \end{minted}
        \caption{Procedimiento \mintinline{pascal}{Merge} del ejercicio~\ref{ej:6}.}
        \label{code:ej6_Merge}
    \end{listing}

    Llamaremos $T_s(k)$ al tiempo que tarda el procedimiento \verb|Sort| cuando actúa sobre un segmento del vector con $k$ entradas. Suponemos que el tiempo que (en media) tarda cada iteración del bucle interno que hay en \verb|Sort| es la unidad (por definición). Es evidente que ese bucle tiene $\dfrac{k(k-1)}{2}$ iteraciones, luego:
    \[
        T_s(k) = \dfrac{k(k-1)}{2} = \dfrac{1}{2}\cdot k^2 - \dfrac{1}{2}\cdot k
    \]

    El tiempo que tarda la versión secuencial sobre $2n$ elementos (llamaremos $S$ a dicho tiempo) será evidentemente $T_s(2n)$, luego:
    \[
        S = T_s(2n) = \dfrac{1}{2}\cdot (2n)^2 - \dfrac{1}{2}\cdot 2n = 2n^2 - n
    \]

    Con estas definiciones, calcular el tiempo que tardará la versión paralela, en los dos siguientes casos. Para esto, hay que suponer que cuando el procedimiento \verb|Merge| actúa sobre un vector con $p$ entradas, tarda $p$ unidades de tiempo en ello, lo cual es razonable teniendo en cuenta que en esas circunstancias \verb|Merge| copia $p$ valores desde \verb|a| hacia \verb|b|. Si llamamos a este tiempo $T_m(p)$, podemos escribir $T_m(p) = p$. Escribe también una comparación cualitativa de los tres tiempos ($S$, $P_1$ y $P_2$). 
    \begin{enumerate}
        \item Las dos instancias concurrentes de \verb|Sort| se ejecutan en el mismo procesador (llamamos $P_1$ al tiempo que tarda).
        
        En este caso, tenemos que hay ganancia de tiempo, ya que las dos instancias de \verb|Sort| no pueden ejecutarse simultáneamente. Por tanto, tenemos que:
        \begin{equation*}
            P_1 = 2\cdot T_s(n) + T_m(2n) = 2\cdot \left( \dfrac{1}{2}\cdot n^2 - \dfrac{1}{2}\cdot n \right) + 2n = n^2 - n + 2n = n^2 + n
        \end{equation*}
        \item Cada instancia de \verb|Sort| se ejecuta en un procesador distinto (lo llamamos $P_2$).
        
        Al poder ejecutarse ahora de forma simultánea, tenemos que:
        \begin{equation*}
            P_2 = \max\{ T_s(n), T_s(n) \} + T_m(2n) = \left( \dfrac{1}{2}\cdot n^2 - \dfrac{1}{2}\cdot n \right) + 2n = \dfrac{1}{2}\cdot n^2 + \dfrac{3}{2}\cdot n
        \end{equation*}
    \end{enumerate}

    Como podemos ver, en los tres casos la eficiencia es del orden cuadrático. No obstante,
    el coeficiente de $n^2$ es distinto en cada caso, siendo el mayor en la versión secuencial. Deducimos por tanto que las versiones concurrentes son más eficientes que la secuencial, y estas mejoras son significativas para valores de $n$ grandes.    

\end{ejercicio}

\begin{ejercicio} \label{ej:7}
    Supongamos que tenemos un programa con tres matrices (\verb|a|, \verb|b| y \verb|c|) de valores flotantes declaradas
    como variables globales. La multiplicación secuencial de \verb|a| y \verb|b| (almacenando el resultado en \verb|c|)
    se puede hacer mediante un procedimiento \verb|MultiplicacionSec| declarado como aparece aquí:
    \begin{minted}{pascal}
        var a, b, c : array[1..3,1..3] of real ;
        procedure MultiplicacionSec()
            var i,j,k : integer ;
            begin
                for i := 1 to 3 do
                    for j := 1 to 3 do begin
                        c[i,j] := 0 ;
                        for k := 1 to 3 do
                            c[i,j] := c[i,j] + a[i,k]*b[k,j] ;
                    end
            end
    \end{minted}
    Escribir un programa con el mismo fin, pero que use 3 procesos concurrentes. Suponer que
    los elementos de las matrices \verb|a| y \verb|b| se pueden leer simultáneamente, así como que elementos distintos de \verb|c| pueden escribirse simultáneamente.\\

    Esta solución puede llevarse a cabo de dos formas (tal y como vimos en la asigantura de AC). Crearemos tres procesos, y cada uno puede calcular una fila de la matriz \verb|c| (Código Fuente~\ref{code:ej7_fila}) o bien cada uno puede calcular una columna de la matriz \verb|c| (Código Fuente~\ref{code:ej7_columna}). No obstante, en la asignatura de AC se vio que, por norma general, es más eficiente calcular por filas que por columnas, ya que en el primer caso se accede a la matriz \verb|a| de forma secuencial, aprovechando así la localidad espacial y provocando un número menor de fallos de caché.
    \begin{figure}
        \centering
        \begin{subfigure}{0.45\textwidth}
            \centering
            \begin{minted}{pascal}
                var a, b, c : array[1..3,1..3] of real ;
                process CalcularFila[ i : 1..3] ;
                    var j, k : integer ;
                    begin
                        for j := 1 to 3 do begin
                            c[i,j] := 0 ;
                            for k := 1 to 3 do
                                c[i,j] := c[i,j] + a[i,k]*b[k,j] ;
                        end
                    end
                end
            \end{minted}
            \caption{Procesos concurrentes para calcular por filas.}
            \label{code:ej7_fila}
        \end{subfigure}\hfill
        \begin{subfigure}{0.45\textwidth}
            \centering
            \begin{minted}{pascal}
                var a, b, c : array[1..3,1..3] of real ;
                process CalcularColumna[j : 1..3] ;
                    var i, k : integer ;
                    begin
                        for i := 1 to 3 do begin
                            c[i,j] := 0 ;
                            for k := 1 to 3 do
                                c[i,j] := c[i,j] + a[i,k]*b[k,j] ;
                        end
                    end
                end
            \end{minted}
            \caption{Procesos concurrentes para calcular por columnas.}
            \label{code:ej7_columna}
        \end{subfigure}
        \caption{Códigos para el ejercicio~\ref{ej:7}.}
    \end{figure}
    % // TODO: Se lanzan solos?
\end{ejercicio}

\begin{ejercicio}\label{ej:8}
    Un trozo de programa ejecuta nueve rutinas o actividades (\verb|P1|, \verb|P2|, \ldots, \verb|P9|), repetidas veces,
    de forma concurrentemente con \verb|cobegin-coend| (ver trozo de código de la figura~\ref{code:ej8_enunciado}), pero que requieren
    sincronizarse según determinado grafo (ver la figura~\ref{fig:ej8_grafo}).
    \begin{figure}
        \centering
        \begin{subfigure}{0.45\textwidth}
            \centering
            \begin{minted}{pascal}
                while true do
                cobegin
                    P1 ; P2 ; P3 ;
                    P4 ; P5 ; P6 ;
                    P7 ; P8 ; P9 ;
                coend
            \end{minted}
            \caption{Código del ejercicio~\ref{ej:8}.}
            \label{code:ej8_enunciado}
        \end{subfigure} \hfill
        \begin{subfigure}{0.45\textwidth}
            \centering
            \begin{tikzpicture}[
                node/.style={circle, draw, minimum size=0.5cm},
                edge/.style={-Stealth}
                ]
    
                % Nodos
                \node[node] (P1) {P1};
                \node[node, right =of P1] (P2) {P2};
                \node[node, right =of P2] (P3) {P3};
                \node[node, below =of P1] (P4) {P4};
                \node[node, right =of P4] (P5) {P5};
                \node[node, right =of P5] (P6) {P6};
                \node[node, below =of P4] (P7) {P7};
                \node[node, right =of P7] (P8) {P8};
                \node[node, right =of P8] (P9) {P9};

                % Aristas
                \draw[edge] (P1) -- (P4);
                \draw[edge] (P1) -- (P2);
                \draw[edge] (P2) -- (P3);
                \draw[edge] (P2) -- (P5);
                \draw[edge] (P3) -- (P6);
                \draw[edge] (P4) -- (P5);
                \draw[edge] (P4) -- (P7);
                \draw[edge] (P5) -- (P6);
                \draw[edge] (P5) -- (P8);
                \draw[edge] (P6) -- (P9);
                \draw[edge] (P7) -- (P8);
                \draw[edge] (P8) -- (P9);
            \end{tikzpicture}
            \caption{DAG del ejercicio~\ref{ej:8}.}
            \label{fig:ej8_grafo}
        \end{subfigure}
        \caption{Figuras del ejercicio~\ref{ej:8}.}
    \end{figure}

    Supón que queremos realizar la sincronización indicada en el grafo, usando para ello llamadas
    desde cada rutina a dos procedimientos (\verb|EsperarPor| y \verb|Acabar|). Se dan los siguientes hechos:
    \begin{itemize}
        \item El procedimiento \verb|EsperarPor(i)| es llamado por una rutina cualquiera (la número $k$) para esperar a que termine la rutina número $i$, usando espera ocupada. Por tanto, se usa por la rutina $k$ al inicio para esperar la terminación de las otras rutinas que corresponda según el grafo.
        \item El procedimiento \verb|Acabar(i)| es llamado por la rutina número $i$, al final de la misma, para indicar que dicha rutina ya ha finalizado.
        \item Ambos procedimientos pueden acceder a variables globales en memoria compartida.
        \item Las rutinas se sincronizan única y exclusivamente mediante llamadas a estos procedimientos, siendo la implementación de los mismos completamente transparente para las rutinas.
    \end{itemize}
    Escribe una implementación de \verb|EsperarPor| y \verb|Acabar| (junto con la declaración e inicialización de las variables compartidas necesarias) que cumpla con los requisitos dados.\\

    Usaremos para ello un vector de variables booleanas \verb|finalizado|, donde \verb|finalizado[i]| indica si la rutina $i$ ha finalizado o no. Inicialmente estará inicializado a \verb|false|, puesto que ningún proceso ha finalizado. Es decir:
    \begin{minted}{pascal}
        var finalizado : array[1..9] of boolean := (false,...,false) ;
    \end{minted}

    El procedimiento \verb|EsperarPor| se implementa de la siguiente forma:
    \begin{minted}{pascal}
        procedure EsperarPor( i : integer ) ;
        begin
            while not finalizado[i] do
            begin
                { no hacer nada }
            end;
        end.
    \end{minted}
    % // TODO: ¿Cuando es procedure y cuándo es process?

    Respecto a la función \verb|Acabar|, podríamos pensar que tan solo se necesita cambiar el valor de la variable \verb|finalizado[i]| a \verb|true|. No obstante, hemos de tener en cuenta que este programa de ejecuta de forma repetida, por lo que si no se reinicia el vector \verb|finalizado| al final de cada ejecución, en la siguiente ejecución ya no habrá la sincronización necesaria. Por tanto, tras el fin del proceso \verb|P9|, se reinicia el vector \verb|finalizado| a \verb|false|.
    \begin{minted}{pascal}
        procedure Acabar( i : integer ) ;
        begin
            finalizado[i] := true ;
            if i >= 9 then
                finalizado := (false,...,false) ;
        end
    \end{minted}
    
\end{ejercicio}


\begin{ejercicio}
    En el ejercicio~\ref{ej:8} los procesos \verb|P1|, \verb|P2|, \ldots, \verb|P9| se ponen en marcha usando \verb|cobegin-coend|.
    Escribe un programa equivalente, que ponga en marcha todos los procesos, pero que use declaración
    estática de procesos, usando un vector de procesos \verb|P|, con índices desde 1 hasta 9, ambos incluidos. El proceso \verb|P[n]| contiene una secuencia de instrucciones desconocida, que llamamos \verb|S_n|, y además debe incluir las llamadas necesarias a \verb|Acabar| y \verb|EsperarPor| (con la misma implementación que antes) para lograr la sincronización adecuada. Se incluye aquí una plantilla:
    \begin{minted}{pascal}
        Process P[ n : 1..9 ]
        begin
            ..... { esperar (si es necesario) a los procesos que corresponda }
            S_n ; { sentencias especificas de este proceso (desconocidas) }
            ..... { senalar que hemos terminado }
        end
    \end{minted}

    En primer lugar, hemos de especificar, en función de $n$, a qué procesos ha de esperar cada uno, lo cual se hará mediante una matriz compartida. Tenemos que:
    \begin{minted}{pascal}
        var espera : array[1..9,1..2] of integer := (
            (-1, -1),   { P1 }
            (1, -1),    { P2 }
            (2, -1),    { P3 }
            (1, -1),    { P4 }
            (2, 4),     { P5 }
            (3, 5),     { P6 }
            (4, -1),    { P7 }
            (5, 7),     { P8 }
            (6, 8)      { P9 }
        ) ;
        Process P[ n : 1..9 ]
        begin
            for i := 1 to 2 do
                if espera[n,i] <> -1 then { != -1 }
                    EsperarPor( espera[n,i] ) ;
            
            S_n ;
            Acabar( n ) ;
        end
    \end{minted}
    % // TODO: Sintaxis en Pascal? Relevante?
\end{ejercicio}

\begin{ejercicio}
    Para los siguientes fragmentos de código, obtener la \emph{poscondición} adecuada para convertirlo en un triple demostrable con la Lógica de Programas:
    \begin{enumerate}
        \item $\{i < 10\} \quad i = 2 \ast i + 1 \quad \{ \}$\\
            Obtenemos la poscondición de este triple razonando matemáticamente:
            \begin{align*}
                i &< 10 \\
                2\ast i &< 20 \\
                i' = 2\ast i + 1 &< 21
            \end{align*}
            donde hemos notado por $i'$ al nuevo valor que adopta la variable $i$.\\

            Por tanto, la poscondición del triple es: $i < 21$.

            Pasamos ahora a demostrar el triple siguiente:
            \begin{equation*}
               \{i < 10\} \quad i = 2 \ast i + 1 \quad \{i<21\}
            \end{equation*}
            
            Usando el axioma de asignación, tenemos que:
            \begin{equation*}
                \{i<21\}^i_{2\ast i + 1}\quad i=2\ast i + 1\quad \{i<21\}
            \end{equation*}

            No obstante, de la definición de Sustitución Textual, tenemos que:
            \begin{equation*}
                \{i<21\}^i_{2\ast i + 1} \equiv \{2\ast i + 1 < 21\} \equiv \{i < 10\} 
            \end{equation*}
            
            Uniendo ambas ecuaciones, obtenemos que el triple es cierto.
            \begin{equation*}
                \{i<10\}\quad i=2\ast i + 1\quad \{i<21\} \\
            \end{equation*}

        \item $\{i > 0\} \quad i = i - 1; \quad \{ \}$\\
            Obtenemos la poscondición de este triple razonando matemáticamente:
            \begin{align*}
                i &> 0 \\
                i' = i - 1 &> -1
            \end{align*}
            donde hemos notado por $i'$ al nuevo valor que adopta la variable $i$.\\

            Por tanto, la poscondición del triple es: $i > -1$.

            Pasamos ahora a demostrar el triple siguiente:
            \begin{equation*}
                \{i > 0\} \quad i = i - 1; \quad \{i > -1\}
            \end{equation*}

            Usando el axioma de asignación, tenemos que:
            \begin{equation*}
                \{i > -1\}^i_{i-1}\quad i=i-1\quad \{i>-1\}
            \end{equation*}

            No obstante, de la definición de Sustitución Textual, tenemos que:
            \begin{equation*}
                \{i > -1\}^i_{i-1} \equiv \{i-1 > -1\} \equiv \{i > 0\}
            \end{equation*}

            Uniendo ambas ecuaciones, obtenemos que el triple es cierto.
            \begin{equation*}
                \{i > 0\}\quad i = i - 1\quad \{i > -1\}
            \end{equation*}

        \item $\{i > j\} \quad i = i + 1;~j = j + 1 \quad \{ \}$\\
            De forma matemática y notando por $i'$ y $j'$ a las modificaciones de $i$ y $j$, respectivamente:
            \begin{align*}
                i &> j \\
                i' = i+1 &> j+1=j'\\
                i' &> j'
            \end{align*}

            Por tanto, la poscondición del triple es: $i > j$. Pasamos a demostrar el triple:
            \begin{equation*}
                \{i > j\} \quad i = i + 1;~j = j + 1 \quad \{i>j\}
            \end{equation*}

            Usando la regla de la composición, tenemos que:
            \begin{equation*}
                \dfrac{\{P\}S_1\{Q\}, \{Q\}S_2\{R\}}{\{P\}S_1;S_2\{R\}} \\
            \end{equation*}

            Identificando $Q$ con $i>j+1$, tenemos que bastará con probar los triples:
            \begin{enumerate}
                \item $\{i>j\}\ i=i+1\ \{i>j+1\}$
                
                Mediante el axioma de asignación, tenemos que:
                \begin{equation*}
                    \{i>j\} \equiv \{i+1>j+1\} \equiv \{i>j+1\}^i_{i+1}\ i=i+1\ \{i>j+1\}
                \end{equation*}

                \item $\{i>j+1\}\ j=j+1\ \{i>j\}$
                
                Mediante el axioma de asignación, tenemos que:
                \begin{equation*}
                    \{i>j+1\} \equiv \{i>j\}^j_{j+1}\ j=j+1\ \{i>j\}
                \end{equation*}
            \end{enumerate}

            Como ambos son ciertos, el triple que queríamos demostrar también lo es gracias a la regla de composición.
        \item $\{\text{falso}\} \quad a = a + 7; \quad \{ \}$\\
            En este caso, partimos de un estado del programa inalcanzable, por lo que en la poscondición podemos poner cualquier estado del programa, es decir, $\{\text{verdad}\}$.
        \item $\{\text{verdad}\} \quad i = 3;~j = 2 \ast i \quad \{ \}$\\
            Como partimos de cualquier estado del programa y sólo se realizan asignaciones, es fácil intuir cuál será la poscondición:
            \begin{gather*}
                i = 3 \\
                j = 2\ast i = 2\ast 3 = 6
            \end{gather*}
            Pasamos a demostrar el triple
            \begin{equation*}
                \{\text{verdad}\} \quad i = 3;~j = 2 \ast i \quad \{i=3\ \land\ j=6\}
            \end{equation*}
            Usando la regla de composición, nos será suficiente probar los triples:
            \begin{equation*}
                \{\text{verdad}\}\ i=3\ \{i=3\} \qquad \{i=3\}\ j=2\ast i\ \{i=3\ \land\ j=6\}
            \end{equation*}
            \begin{enumerate}
                \item Para el primer triple, usamos el axioma de asignación:
                \begin{equation*}
                    \{\text{verdad}\} \equiv \{3=3\} \equiv \{i=3\}^i_3\ i=3\ \{i=3\}
                \end{equation*}
                \item Para el segundo, volvemos a usar el axioma de asignación:
                \begin{equation*}
                    \{i=3\ \land\ j=6\}^j_{2\ast i}\ j=2\ast i\ \{i=3\ \land\ j=6\}
                \end{equation*}

                No obstante, de la definición de Sustitución Textual, tenemos que:
                \begin{multline*}
                    \{i=3\ \land\ j=6\}^j_{2\ast i} \equiv \{i=3\ \land\ 2\ast i=6\} \equiv \{i=3\ \land 2\ast 3=6\} \equiv \\ \equiv \{i=3\ \land\ 6=6\} \equiv \{i=3\}
                \end{multline*}

                Uniendo ambas ecuaciones, obtenemos que el triple es cierto.
            \end{enumerate}
            Ambos triples son ciertos, luego por la regla de la composición tenemos demostrado nuestro triple.
        \item $\{\text{verdad}\} \quad c = a + b;~c = \nicefrac{c}{2} \quad \{ \}$\\
            Notando por $c'$ al nuevo valor de $c$, tenemos que:
            \begin{align*}
                c &= a + b \\
                c' = \nicefrac{c}{2} &= \frac{a+b}{2}
            \end{align*}
            Tratamos por tanto de probar el siguiente triple
            \begin{equation*}
                \{\text{verdad}\} \quad c = a + b;~c = \nicefrac{c}{2} \quad \left\{c=\frac{a+b}{2}\right\}
            \end{equation*}
            Usando la regla de la composición, basta con probar los triples:
            \begin{equation*}
                \{\text{verdad}\}\ c=a+b;\ \{c=a+b\} \qquad \{c=a+b\}\ c=\nicefrac{c}{2};\ \left\{c=\frac{a+b}{2}\right\}
            \end{equation*}
            \begin{enumerate}
                \item Para el primero, usamos el axioma de asignación:
                    \begin{equation*}
                        \{\text{verdad}\}\equiv\{a+b=a+b\}\equiv \{c=a+b\}^c_{a+b}\ c=a+b\ \{c=a+b\}
                    \end{equation*}
                \item Para la segunda, también usamos el axioma de asignación:
                    \begin{equation*}
                        \{c=a+b\}\equiv \left\{\frac{c}{2} = \frac{a+b}{2}\right\} \equiv \left\{c=\frac{a+b}{2}\right\}^c_{\nicefrac{c}{2}}\ c=\nicefrac{c}{2}\ \left\{c=\frac{a+b}{2}\right\}
                    \end{equation*}
            \end{enumerate}
            Usando la relga de composición, tenemos demostrado nuestro triple.
    \end{enumerate}
\end{ejercicio}

\begin{ejercicio}
    ¿Cuáles de los siguientes triples no son demostrables con la Lógica de Programas?
    (Considerando que $i,x,a \in \mathbb{Z}$)
    \begin{enumerate}
        \item $\{i > 0\} \quad i = i - 1; \quad \{i \geq 0\}$\\
            El siguiente triple sabemos que es cierto:
            \begin{equation*}
                \{i>0\}\ i=i-1\ \{i>-1\} 
            \end{equation*}
            $\{i>-1\}\rightarrow \{i\geq 0\}$, luego es cierto por la primera regla de la consecuencia.
        \item $\{x \geq 7\} \quad x = x + 3; \quad \{x \geq 9\}$\\
            El siguiente triple sabemos que es cierto:    
            \begin{equation*}
                \{x\geq 7\}\ x=x+3\ \{x\geq 10\}
            \end{equation*}
            $\{x\geq 10\}\rightarrow\{x\geq 9\}$, luego es cierto por la primera regla de la consecuencia.
        \item $\{i < 9\} \quad i = 2 \ast i + 1; \quad \{ i \leq 20\}$\\
            El siguiente triple sabemos que es cierto:
            \begin{equation*}
                \{i<9\}\ i=2\ast i+1\ \{i<19\}
            \end{equation*}
            $\{i<19\}\rightarrow\{i\leq 20\}$, luego es cierto por la primera regla de la consecuencia.
        \item $\{a > 0\} \quad a = a - 7; \quad \{a > -6\}$\\
            \begin{equation*}
                \{a>0\}\ a=a-7\ \{a>-7\}
            \end{equation*}
            Pero $\{a>-7\}\not\rightarrow\{a>-6\}$, luego este triple no es demostrable.
    \end{enumerate}
\end{ejercicio}

\begin{ejercicio}
    Si el triple $\{P\} C \{Q\}$ es demostrable, indicar por qué los siguientes triples también lo son (o no se pueden demostrar y por qué):
    \begin{enumerate}
        \item $\{P\} C \{Q \lor P\}$\\
           Es demostrable, ya que $\{Q\}\rightarrow\{Q \lor P\}$ y por la primera regla de la consecuencia, tomando $R=Q \lor P$:
            \begin{equation*}
                \dfrac{\{P\}C\{Q\}, \{Q\}\rightarrow\{R\}}{\{P\}C\{R\}}
            \end{equation*}
            Tenemos que se debilita la poscondición.

        \item $\{P \land D\} C \{Q\}$\\
            Es demostrable, ya que $\{P \land D\}\rightarrow\{P\}$ y por la segunda regla de la consecuencia, tomando $R = P \land D$:
            \begin{equation*}
                \dfrac{\{P\}\rightarrow\{R\},\{R\}C\{Q\}}{\{P\}C\{Q\}}
            \end{equation*}
            Tenemos que se fortalece la precondición.
        \item $\{P \lor D\} C \{Q\}$\\
            No es demostrable, porque se debilita la precondición.

        \item $\{P\} C \{Q \lor D\}$\\
            Al igual que hemos hecho en el apartado 1, es demostrable ya que ${\{Q\}\rightarrow\{Q \lor D\}}$ y usando la primera regla de la consecuencia.
        \item $\{P\} C \{Q \land P\}$\\
            No podemos demostrarlo, ya que se fortalece la poscondición.
    \end{enumerate}
\end{ejercicio}

\begin{ejercicio}
    Si el triple $\{P\} C \{Q\}$ es demostrable, ¿cuál de los siguientes triples no se puede demostrar?
    \begin{enumerate}
        \item $\{P \land D\} C \{Q\}$\\
            Sabemos que $\{P \land D\}\rightarrow\{P\}$, luego puede demostrarse por la segunda regla de la consecuencia (se fortalece la precondición).
        \item $\{P \lor D\} C \{Q\}$\\
            No puede demostrarse, porque se debilita la precondición.
        \item $\{P\} C \{Q \lor D\}$\\
            Puede demostrarse mediante la primera regla de la consecuencia, ya que se tiene que ${\{Q\}\rightarrow\{Q \lor D\}}$.

        \item $\{P\} C \{Q \lor P\}$\\
        Puede demostrarse mediante la primera regla de la consecuencia, ya que se tiene que ${\{Q\}\rightarrow\{Q \lor P\}}$.
    \end{enumerate}
\end{ejercicio}

\begin{ejercicio}
    Dado el siguiente programa, obtener:
    \begin{minted}{pascal}
        int x = 5, y = 2;
        cobegin
            < x = x + y >;
            < y = x * y >;
        coend
    \end{minted}
    \begin{enumerate}
        \item Valores finales de $x$ e $y$.
        Tenemos dos posibles trazas de ejecución:
        \begin{enumerate}
            \item Primero se ejecuta la primera instrucción, por lo que obtendríamos $x = 7$ y $y = 14$.
            \item Primero se ejecuta la segunda instrucción, por lo que obtendríamos $x = 15$ y $y = 10$.
        \end{enumerate}
        \item Valores finales de $x$ e $y$ si quitamos los símbolos \verb|< >| de instrucción atómica.
        
        Encontramos cada uno de los dos estados anteriores, además de $x = 7$ y $y= 10$.
    \end{enumerate}
\end{ejercicio}

\begin{ejercicio}
    Comprobar si la demostración del siguiente triple interfiere con los teoremas siguientes:
    \[
        \{x \geq 2\} \quad < x = x - 2 > \quad \{x \geq 0\}
    \]
    Es decir, queremos comprobar si $R\equiv <x = x-2>$ con $pre(R) = \{x\geq 2\}$ interfiere con los triples siguientes:
    \begin{enumerate}
        \item $\{x \geq 0\} \quad < x = x + 3 > \quad \{x \geq 3\}$\\
            Comprobamos en primer lugar su interferencia con la precondición:
            \begin{equation*}
                \{x\geq 0 \land x \geq 2\}\ <x=x-2>\ \{x\geq 0\}
            \end{equation*}
            Este triple es correcto por la segunda regla de la consecuencia, luego no interfiere con la precondición.\\

            Comprobemos ahora su interferencia con la poscondición:
            \begin{equation*}
                \{x\geq 3 \land x \geq 2\}\ <x=x-2>\ \{x\geq 1\}
            \end{equation*}
            En este caso, $\{x\geq 1\}\not\rightarrow\{x\geq 3\}$, luego este triple no es demostrable y $R$ interfiere con la poscondición del triple en cuestión.

        \item $\{x \geq 0\} \quad < x = x + 3 > \quad \{x \geq 0\}$\\
            Comprobamos en primer lugar su interferencia con la precondición:
            \begin{equation*}
                \{x\geq 0 \land x \geq 2\}\ <x=x-2>\ \{x\geq 0\}
            \end{equation*}
            Este triple es correcto por la segunda regla de la consecuencia, luego no interfiere con la precondición.
            Además, como la precondición y la poscondición son iguales, $R$ tampoco interfiere con la poscondición, luego no interfiere con este triple.
        \item $\{x \geq 7\} \quad < x = x + 3 > \quad \{x \geq 10\}$\\
            Comprobamos en primer lugar su interferencia con la precondición:
            \begin{equation*}
                \{x\geq 7 \land x \geq 2\}\ <x=x-2>\ \{x\geq 5\}
            \end{equation*}
            No obstante, como $\{x\geq 5\}\not\rightarrow\{x\geq 7\}$, $R$ interfiere con la precondición de este triple.
        \item $\{y \geq 0\} \quad < y = y + 3 > \quad \{y \geq 3\}$\\
            $R$ no interfiere con este triple, ya que son variables disjuntas.
        \item $\{x \text{ es impar}\} \quad < y = x + 1 > \quad \{y \text{ es par}\}$\\
            Comprobamos en primer lugar su interferencia con la precondición:
            \begin{equation*}
                \{x \text{ es impar } \land x\geq 2\}\ <x=x-2>\ \{x \text{ es impar } \land x \geq 0\}
            \end{equation*}
            Por la $2^a$ regla de la consecuencia, como $\{x \text{ es impar } \land x \geq 0\} \rightarrow \{x \text{ es impar}\}$, tenemos que es correcto y $R$ no interfiere con la precondición.

            Comprobamos ahora su interferencia con la poscondición:
            \begin{equation*}
                \{y \text{ es par } \land x \geq 0\}\ <x=x-2>\ \{y \text{ es par } \land x \geq -2\}
            \end{equation*}
            Por la $2^a$ regla de la consecuencia, como $\{y \text{ es par } \land x \geq -2\} \rightarrow \{y \text{ es par}\}$, tenemos que es correcto y $R$ no interfiere con la poscondición. Por tanto, $R$ no interfiere con este triple.
            
    \end{enumerate}
\end{ejercicio}

\begin{ejercicio}
    Dado el siguiente triple:
    \begin{gather*}
        \{x = 0\} \\
        \text{cobegin} \\
        <x = x + a> || <x = x + b> || <x = x + c> \\
        \text{coend} \\
        \{x = a + b + c\}
    \end{gather*}
    
    Demostrarlo utilizando la lógica de asertos para cada una de las tres instrucciones atómicas y después que se llega a la poscondición final $x = a + b + c$ utilizando para ello la regla \emph{de la composición concurrente} de instrucciones atómicas.\\

Inicialmente, demostraremos los $3$ siguientes triples, uno por cada instrucción atómica. Hemos de notar que, en cada uno de ellos, como no sabemos en qué orden se ejecutan, tenemos que incluir en las precondiciones y las poscondiciones todas las posibilidades.
\begin{enumerate}
    \item El correspondiente a la primera instrucción atómica:
    \begin{gather*}
        \{x=0 \lor x=b \lor x=c \lor x=b+c\}\ <x=x+a>\ \\ \{x=a \lor x=a+b \lor x=a+c \lor x=a+b+c\}
    \end{gather*}

    Mediante el axioma de asignación, tenemos que:
    \begin{gather*}
        \{x=a \lor x=a+b \lor x=a+c \lor x=a+b+c\}_{x+a}^a <x=x+a>\\ \{x=a \lor x=a+b \lor x=a+c \lor x=a+b+c\}
    \end{gather*}

    No obstante, de la definición de Sustitución Textual, tenemos:
    \begin{gather*}
        \{x=a \lor x=a+b \lor x=a+c \lor x=a+b+c\}_{x+a}^a \equiv \\
        \equiv \{x+a=a \lor x+a=a+b \lor x+a=a+c \lor x+a=a+b+c\} \equiv \\
        \equiv \{x=0 \lor x=b \lor x=c \lor x=b+c\}
    \end{gather*}

    Por tanto, el triple en cuestión es cierto.

    \item El correspondiente a la segunda instrucción atómica:
    \begin{gather*}
        \{x=0 \lor x=a \lor x=c \lor x=a+c\}\ <x=x+b>\ \\ \{x=b \lor x=a+b \lor x=b+c \lor x=a+b+c\}
    \end{gather*}

    Es cierto, y su demostración es análoga al primer caso.
    \item El correspondiente a la tercera instrucción atómica:
    \begin{gather*}
        \{x=0 \lor x=b \lor x=a \lor x=a+b\}\ <x=x+c>\ \\ \{x=c \lor x=a+c \lor x=b+c \lor x=a+b+c\}
    \end{gather*}
    Es cierto, y su demostración es análoga al primer caso.
\end{enumerate}

Seguidamente, tenemos que ver que dichos $3$ triples están libres de interferencias. Para ello, hemos de probar $12$ triples, ya que hay $3$ instrucciones atómicas, cada una de ellas con $2$ asertos, por lo que por cada instrucción hemos de comprobar $4$ asertos:
\begin{align*}
    &NI(x=0 \lor x=a \lor x=c \lor x=a+c, <x=x+a>)\\
    &NI(x=b \lor x=a+b \lor x=b+c \lor x=a+b+c, <x=x+a>)\\
    &NI(x=0 \lor x=b \lor x=a \lor x=a+b, <x=x+a>)\\
    &NI(x=c \lor x=a+c \lor x=b+c \lor x=a+b+c, <x=x+a>)\\ \\
    &NI(x=0 \lor x=b \lor x=c \lor x=b+c, <x=x+b>)\\
    &NI(x=a \lor x=a+b \lor x=a+c \lor x=a+b+c, <x=x+b>)\\
    &NI(x=0 \lor x=b \lor x=a \lor x=a+b, <x=x+b>)\\
    &NI(x=c \lor x=a+c \lor x=b+c \lor x=a+b+c, <x=x+b>)\\ \\
    &NI(x=0 \lor x=b \lor x=c \lor x=b+c, <x=x+c>)\\
    &NI(x=a \lor x=a+b \lor x=a+c \lor x=a+b+c, <x=x+c>)\\
    &NI(x=0 \lor x=a \lor x=c \lor x=a+c, <x=x+c>)\\
    &NI(x=b \lor x=a+b \lor x=b+c \lor x=a+b+c, <x=x+c>)\\
\end{align*}

Demostremos ahora el primero, ya que el resto son idénticos.
\begin{align*}
    NI&(x=0 \lor x=a \lor x=c \lor x=a+c, <x=x+a>)\equiv \\
    &\equiv \{(x=0 \lor x=a \lor x=c \lor x=a+c) \land (x=0 \lor x=b \lor x=c \lor x=b+c)\}\\&\qquad <x=x+a> \{x=0 \lor x=a \lor x=c \lor x=a+c\}\equiv \\
    &\equiv \{x=0 \lor x=c\} <x=x+a> \{x=0 \lor x=a \lor x=c \lor x=a+c\}
\end{align*}
Este triple efectivamente es cierto, lo cual se puede demostrar empleando en primer lugar el Axioma de Sustitución y, posteriormente, la primera regla de la consecuencia.\\

Por tanto, y tras aplicar la Regla de la Composición Concurrente, tenemos de forma directa que:
    \begin{gather*}
        \{x = 0\} \\
        \text{cobegin} \\
        <x = x + a> || <x = x + b> || <x = x + c> \\
        \text{coend} \\
        \{x = a + b + c\}
    \end{gather*}
\end{ejercicio}

\begin{ejercicio}
    El siguiente triple:
    \begin{gather*}
        \{x=0 \land y = 0 \land z = 0\} \\
        <x = z + a> || <y = x+b> \\
        \{(x = a) \land (y = b \lor y = a +b) \land z = 0\}
    \end{gather*}
    \begin{enumerate}[label=(\alph*)]
        \item Es indemostrable salvo que se cumpla siempre que $a=0$.
        \item El triple anterior es demostrable para cualquier valor de las variables $a$ o $b$.
        \item Es indemostrable salvo que se cumpla siempre que $b = 0$.
        \item Es indemostrable salvo que se cumpla siempre que $a = 0 \land b = 0$.
    \end{enumerate}

    Veamos si podemos demostrarlo. Para ello, notamos cada instrucción atómica de la siguiente forma:
    \begin{align*}
        S_1 &= <x = z + a> \\
        S_2 &= <y = x + b>
    \end{align*}

    Veamos cuál ha de ser la precondición de cada instrucción atómica:
    \begin{align*}
        P_1 &= x=0 \land (y=0 \lor y=b) \land z=0 \\
        P_2 &= (x=0\lor x=a) \land y=0 \land z=0
    \end{align*}

    Veamos cuál ha de ser la poscondición de cada instrucción atómica:
    \begin{align*}
        Q_1 &= x=a\land (y=0 \lor y=b)\land z=0 \\
        Q_2 &= [(x=0\land y=b) \lor (x=a\land y=a+b)] \land z=0
    \end{align*}

    Vemos por tanto que ambos triples son ciertos:
    \begin{enumerate}
        \item $\{P_1\} S_1 \{Q_1\}$
        
        De la definición de Sustitución Textual, tenemos que:
        \begin{multline*}
            \{x=a \land (y=0 \lor y=b) \land z = 0\}_{z+a}^x \equiv \{z+a=a \land (y=0 \lor y=b) \land z = 0\} \equiv\\\equiv \{(y=0 \lor y=b) \land z=0\}
        \end{multline*}

        Por tanto, del Axioma de Asignación, tenemos que:
        \begin{equation*}
            \{(y=0 \lor y=b) \land z = 0\} <x = z + a> \{x=a \land (y=0 \lor y=b) \land z = 0\}
        \end{equation*}

        Finalmente, usando la segunda regla de la consecuencia, como se tiene que $\{P_1\}\rightarrow\{(y=0 \lor y=b) \land z = 0\}$, tenemos que el triple es cierto.
        
        \item $\{P_2\} S_2 \{Q_2\}$
        
        De la definición de Sustitución Textual, tenemos que:
        \begin{multline*}
            \{[(x=0\land y=b) \lor (x=a\land y=a+b)] \land z=0\}_{x+b}^y \equiv \\ \equiv \{[(x=0\land x+b=b) \lor (x=a\land x+b=a+b)] \land z = 0\} \equiv \\ \equiv \{(x=0\lor x=a) \land z = 0\}
        \end{multline*}

        Por tanto, del Axioma de Asignación, tenemos que:
        \begin{multline*}
            \{(x=0\lor x=a) \land z = 0\} <y = x + b> \\ \{[(x=0\land y=b) \lor (x=a\land y=a+b)] \land z=0\}
        \end{multline*}

        Por tanto, usando la segunda regla de la consecuencia, como se tiene que $\{P_2\}\rightarrow\{(x=0\lor x=a) \land z = 0\}$, tenemos que el triple es cierto.
    \end{enumerate}

    Veamos ahora que no interfieren entre sí. Como tenemos dos instrucciones atómicas, cada una con dos asertos, hemos de comprobar $4$ asertos:
    \begin{align*}
        NI(P_2, S_1) & \equiv \{P_1\land P_2\} S_1 \{P_2\} \\
        NI(Q_2, S_1) & \equiv \{P_1\land Q_2\} S_1 \{Q_2\} \\
        NI(P_1, S_2) & \equiv \{P_2\land P_1\} S_2 \{P_1\} \\
        NI(Q_1, S_2) & \equiv \{P_2\land Q_1\} S_2 \{Q_1\}
    \end{align*}

    Veamos por tanto en qué queda cada uno de ellos:
    \begin{align*}
        NI(P_2, S_1) & \equiv \{x=0 \land y=0 \land z=0\} <x=z+a> \{(x=0\lor x=a) \land y=0 \land z=0\} \\
        NI(Q_2, S_1) & \equiv \{x=0 \land y=b \land z=0\} <x=z+a>\\&\hspace{5cm} \{[(x=0\land y=b) \lor (x=a\land y=a+b)] \land z=0\} \\
        NI(P_1, S_2) & \equiv \{x=0 \land y=0 \land z=0\} <y=x+b> \{x=0 \land (y=0\lor y=b) \land z=0\} \\
        NI(Q_1, S_2) & \equiv \{x=a \land y=0 \land z=0\} <y=x+b> \{x=a \land (y=0\lor y=b) \land z=0\}
    \end{align*}

    Intentemos demostrar el segundo. Usando la definición de Sustitución Textual, tenemos que:
    \begin{multline*}
        \{[(x=0\land y=b) \lor (x=a\land y=a+b)] \land z=0\}_{z+a}^x \equiv \\ \equiv \{[(z+a=0\land y=b) \lor (z+a=a\land y=a+b)] \land z=0\} \equiv \\ \equiv \{y=a+b \land z=0\}
    \end{multline*}
    Por tanto, del Axioma de Asignación, tenemos que:
    \begin{multline*}
        \{y=a+b \land z=0\} <x=z+a> \{[(x=0\land y=b) \lor (x=a\land y=a+b)] \land z=0\}
    \end{multline*}
    No obstante, de forma general, tenemos que $\{x=0\land y=b\land z=0\}\not\rightarrow\{y=a+b \land z=0\}$, por lo que $NI(Q_2, S_1)$ no es demostrable. No obstante, si $a=0$, entonces sí que sería demostrable. Supongamos por tanto a partir de ahora que \ul{$a=0$}.

    Intentemos ahora demostrar el cuarto. Usando la definición de Sustitución Textual, tenemos que:
    \begin{multline*}
        \{x=a \land (y=0 \lor y=b) \land z = 0\}_{x+b}^y \equiv \\ \equiv \{x=a \land (x+b=0 \lor x+b=b) \land z = 0\} \equiv \\ \equiv \{x=a \land (x=-b\lor x=0) \land z = 0\} \equiv 
        \{x=0 \land z = 0\}
    \end{multline*}
    donde en la última igualdad hemos usado que \ul{$a=0$}. Por tanto, del Axioma de Asignación, tenemos que:
    \begin{multline*}
        \{x=0 \land z = 0\} <y=x+b> \{x=a \land (y=0 \lor y=b) \land z = 0\}
    \end{multline*}
    Además, como \ul{$a=0$}, tenemos que $\{x=0\land z=0\}\rightarrow\{x=a \land y=0 \land z=0\}$, tenemos que es cierto. Por tanto, $NI(P_1, S_2)$ es demostrable.

    Los otros dos triples son análogamente ciertos, por lo que podemos aplicar la regla de la composición concurrente y llegar al siguiente triple:
    \begin{gather*}
        \{x=0 \land y = 0 \land z = 0\} \\
        <x = z + a> || <y = x+b> \\
        \{x = a \land  y = a +b \land z = 0\}
    \end{gather*}
    Este es el triple que queríamos demostrar, suponiendo que \ul{$a=0$}. Por tanto, la respuesta correcta es la \textbf{a)}, ya que no es demostrable salvo que se cumpla siempre que $a=0$.

    % // TODO: JJ Revisar. Me ha parecido bastante difícil
\end{ejercicio}

\begin{ejercicio}
    Suponer que $\{suma>1\}\ suma=suma+4\ \{suma>5\}$ es demostrable, entonces: ¿cuál de los siguientes triples es también demostrable? (indicar por qué)
    \begin{enumerate}
        \item $\{suma>2\}\ suma=suma+4\ \{suma>5\}$.\\
            Es demostrable, ya que $\{suma>2\}\rightarrow\{suma>1\}$ y podemos aplicar la segunda regla de la consecuencia.
        \item $\{suma\geq 1\}\ suma=suma+4\ \{suma > 5\}$.\\
            No es demostarble, ya que debilita la precondición.
        \item $\{suma>0\}\ suma=suma+4\ \{suma > 5\}$.\\
            No es demostarble, ya que debilita la precondición.
        \item $\{suma>1\}\ suma=suma+4\ \{suma > 6\}$.\\
            No es demostrable, ya que fortalece la poscondición.
    \end{enumerate}
\end{ejercicio}

\begin{ejercicio}
    Suponer que $\{x<y\}\ C_1\ \{u<v\}$ es demostrable, entonces: ¿cuáles de los siguientes triples son también demostrables? (indicar por qué)
    \begin{enumerate}
        \item $\{x\leq y\}\ C_1\ \{u<v\}$.\\
            No es demostrable, ya que debilita la precondición.
        \item $\{x\leq y-2\}\ C_1\ \{u<v\}$.\\
            Es demostrable, ya que $\{x\leq y-2\}\rightarrow\{x+2\leq y\}\rightarrow\{x<y\}$, y mediante la segunda regla de la consecuencia se tiene que es cierto.
        \item $\{x\leq y\}\ C_1\ \{u\leq v\}$.\\
            El triple $\{x<y\}\ C_1\ \{u\leq v\}$ sí que es demostrable ya que relaja la poscondición, pero el triple que se nos dice no es demostrable, ya que también relaja la precondición.
            Como además no tenemos relación entre $x$ y $u$ ni entre $y$ y $v$, no podemos inferir nada.
        \item $\{x<y\}\ C_1\ \{u<v-2\}$.\\
            No es demostrable, ya que fortalecemos la poscondición.
    \end{enumerate}
\end{ejercicio}

\begin{ejercicio}
    Seleccionar el valor correcto de las 2 variables ($x$ e $y$) después de ejecutarse el siguiente programa concurrente:
    \begin{minted}{pascal}
        int x=5, y=2;
        cobegin <x=x+y>; <y=x*y>; <x=x-y>; coend;
    \end{minted}
    \begin{enumerate}[label=(\alph*)]
        \item $x=7$ y $y=14$.
        \item $x=5$ y $y=10$.
        \item $x=-7$ y $y=14$.
        \item $x=-3$ y $y=10$.
    \end{enumerate}

    Numeramos las instrucciones atómicas de la siguiente forma:
    \begin{enumerate}
        \item \verb|<x=x+y>|
        \item \verb|<y=x*y>|
        \item \verb|<x=x-y>|
    \end{enumerate}

    Veamos ahora, en función del orden de ejecución, cuál sería el valor de las variables $x$ e $y$:
    \begin{itemize}
        \item \ul{1, 2, 3}: $x=-7$ y $y=14$.
        \item \ul{1, 3, 2}: $x=5$ y $y=10$.
        \item \ul{3, 1, 2}: $x=5$ y $y=10$.
        \item \ul{2, 1, 3}: $x=5$ y $y=10$.
        \item \ul{2, 3, 1}: $x=5$ y $y=10$.
        \item \ul{3, 2, 1}: $x=9$ y $y=6$.
    \end{itemize}

    Por tanto, las respuestas $b$ y $c$ son correctas.
    % // TODO: Lo ha hecho Arturo. Revisar JJ
\end{ejercicio}

\begin{ejercicio}
    El siguiente código concurrente no puede ser demostrado directamente con la lógica de aserciones (pre y poscondiciones). Elegir la respuesta que explica correctamente la razón de que ocurra esto.
    \begin{minted}{pascal}
        {x=0} cobegin <x=x+a>; <x=x+a> coend; {x=2*a}
        (a es un valor entero positivo)
    \end{minted}
    \begin{enumerate}[label=(\alph*)]
        \item Porque la poscondición que se propone $\{x=2\ast a\}$ es falsa.
        \item Porque falta incluir la posibilidad de que el valor final de $x$ sea también ${\{x=a\}}$.
        \item Porque al aplicar directamente la regla de inferencia de la \textit{composición concurrente} utilizo unas condiciones (pre y post-condiciones) demasiado débiles.
        \item Porque tengo que incluir en los asertos el valor del contador de programa de cada procesador.
    \end{enumerate}

    Notamos cada instrucción atómica de la siguiente forma:
    \begin{equation*}
        S_1 = S_2 = <x=x+a>
    \end{equation*}

    Veamos cuál ha de ser la precondición de cada instrucción atómica:
    \begin{equation*}
        P_1 = P_2 = x=0 \lor x=a
    \end{equation*}

    Veamos cuál ha de ser la poscondición de cada instrucción atómica:
    \begin{equation*}
        Q_1 = Q_2 = x=a \lor x=2a
    \end{equation*}

    Veamos ahora que cada triple es cierto. Como son los mismos, hemos de demostrar:
    \begin{equation*}
        \{x=0 \lor x=a\} <x=x+a> \{x=a \lor x=2a\}
    \end{equation*}
    
    Este se demuestra de forma directa. Además, también hemos de demostrar que no interfieren entre sí. Tenemos que demostrar:
    \begin{align*}
        NI(P_1, S_2) &\equiv \{x=0 \lor x=a\} <x=x+a> \{x=0 \lor x=a\} \\
        NI(Q_1, S_2) &\equiv \{x=a\} <x=x+a> \{x=a \lor x=2a\}
    \end{align*}

    Estos también son ciertos, por lo que podemos aplicar la regla de la composición concurrente y llegar al siguiente triple:
    \begin{equation*}
        \{x=0 \lor x=a\} cobegin <x=x+a>; <x=x+a> coend; \{x=a\lor x=2a\}
    \end{equation*}

    Por la primera regla de la consecuencia, podemos debilitar la precondición, llegando al siguiente triple:
    \begin{equation*}
        \{x=0\} cobegin <x=x+a>; <x=x+a> coend; \{x=a\lor x=2a\}
    \end{equation*}
    
    No obstante, la poscondición no se puede debilitar, por lo que la respuesta correcta es la \textbf{b)}, ya que falta incluir en los asertos el valor final de $x$ sea también $\{x=a\}$.

    % // TODO: Preguntar a Capel
    % {x=0} cobegin <x=x+a>; <x=x+a> coend; {x=2*a}       no se puede demostrar
    % {x=0} cobegin <x=x+a>; <x=x+a> coend; {x=2*a  o   x=a}    es cierto
\end{ejercicio}

\begin{ejercicio}
   Estudiar cuáles son los valores finales de las variables \verb|x| e \verb|y| en el siguiente programa. Insertar los asertos adecuados entre llaves, antes y después de cada sentencia, para poder obtener una traza de demostración del programa, que incluya en su último aserto los valores finales de las variables. 
   \begin{minted}{c++}
       int x = c1;
       int y = c2;
       x = x + y;
       y = x * y;
       x = x - y;
   \end{minted}

   Tenemos que cada triple, por orden, es:
   \begin{gather*}
       \{x = c_1 \land y = c_2\}\\
       x = x + y \\
       \{x = c_1 + c_2 \land y =c_2\} \\
       y = x * y \\
       \{x = c_1 + c_2 \land y = (c_1 + c_2)\cdot c_2 \} \\
       x = x - y \\
       \{x = (c_1 + c_2) - (c_1 + c_2)\cdot c_2 = (c_1 + c_2)\cdot (1-c_2) \land y = (c_1 + c_2)\cdot c_2 \}
   \end{gather*}
\end{ejercicio}

\begin{ejercicio}
    Demostrar que el siguiente triple es cierto:
    \begin{gather*}
        \{x=0\}\\ 
        \texttt{cobegin} \\
        <x=x+1> || <x=x+2> || <x=x+4> \\
        \texttt{coend} \\
        \{x=7\}
    \end{gather*}

    Notamos cada instrucción atómica de la siguiente forma:
    \begin{align*}
        S_1 &= <x=x+1> \\
        S_2 &= <x=x+2> \\
        S_3 &= <x=x+4>
    \end{align*}

    Veamos cuál ha de ser la precondición de cada instrucción atómica:
    \begin{align*}
        P_1 = x=0 \lor x=2 \lor x=4 \lor x=6 \\
        P_2 = x=0 \lor x=1 \lor x=4 \lor x=5 \\
        P_3 = x=0 \lor x=1 \lor x=2 \lor x=3
    \end{align*}

    Veamos cuál ha de ser la poscondición de cada instrucción atómica:
    \begin{align*}
        Q_1 = x=1 \lor x=3 \lor x=5 \lor x=7 \\
        Q_2 = x=2 \lor x=3 \lor x=6 \lor x=7 \\
        Q_3 = x=4 \lor x=5 \lor x=6 \lor x=7
    \end{align*}

    Cada triple es directamente cierto por el Axioma de Asignación. Veamos ahora que no interfieren entre sí. Tenemos que demostrar:
    \begin{align*}
        NI(P_2, S_1) &\equiv \{x=0 \lor x=4\} <x=x+1> \{x=0\lor x=1\lor x=4\lor x=5\} \\
        NI(Q_2, S_1) &\equiv \{x=2 \lor x=6\} <x=x+1> \{x=2\lor x=3\lor x=6\lor x=7\} \\
        NI(P_3, S_1) &\equiv \{x=0 \lor x=2\} <x=x+1> \{x=0\lor x=1\lor x=2\lor x=3\} \\
        NI(Q_3, S_1) &\equiv \{x=4 \lor x=6\} <x=x+1> \{x=4\lor x=5\lor x=6\lor x=7\} \\ \\
        NI(P_1, S_2) &\equiv \{x=0 \lor x=4\} <x=x+2> \{x=0\lor x=2\lor x=4\lor x=6\} \\
        NI(Q_1, S_2) &\equiv \{x=1 \lor x=5\} <x=x+2> \{x=1\lor x=3\lor x=5\lor x=7\} \\
        NI(P_3, S_2) &\equiv \{x=0 \lor x=1\} <x=x+2> \{x=0\lor x=1\lor x=2\lor x=3\} \\
        NI(Q_3, S_2) &\equiv \{x=4 \lor x=5\} <x=x+2> \{x=4\lor x=5\lor x=6\lor x=7\} \\ \\
        NI(P_1, S_3) &\equiv \{x=0 \lor x=2\} <x=x+4> \{x=0\lor x=2\lor x=4\lor x=6\} \\
        NI(Q_1, S_3) &\equiv \{x=1 \lor x=3\} <x=x+4> \{x=1\lor x=3\lor x=5\lor x=7\} \\
        NI(P_2, S_3) &\equiv \{x=0 \lor x=1\} <x=x+4> \{x=0\lor x=1\lor x=4\lor x=5\} \\
        NI(Q_2, S_3) &\equiv \{x=2 \lor x=3\} <x=x+4> \{x=2\lor x=3\lor x=6\lor x=7\}
    \end{align*}

    Todos estos son ciertos, por lo que podemos aplicar la regla de la composición concurrente y llegar al siguiente triple:
    \begin{gather*}
        \{x=0\} \\
        \texttt{cobegin} \\
        <x=x+1> || <x=x+2> || <x=x+4> \\
        \texttt{coend} \\
        \{x=7\}
    \end{gather*}

    % // TODO: JJ Revisar
    % // Este ha sido fácil. Intuición? Cómo verlo directo sin escribir todo?
\end{ejercicio}

\begin{ejercicio}
    Dada la siguiente construcción de composición concurrente \verb|P|:
    \begin{gather*}
        \texttt{cobegin} \\
        <x=x-1>;<x=x+1>;||<y=y-1>;<y=y+1>; \\
        \texttt{coend} \\
    \end{gather*}
    demostrar que se cumple la invarianza de $\{x=y\}$, es decir, que $\{x=y\}\ P\ \{x=y\}$ es un triple cierto.\\

Para ello, comenzamos demostrando los siguientes triples (se ha de mantener la invarianza en el código secuencial)
\begin{gather*}
    \{x=y\}\ x=x-1;x=x+1;\ \{x=y\} \\
    \{x=y\}\ y=y-1;y=y+1;\ \{x=y\} \\
\end{gather*}
\begin{enumerate}
    \item Respecto al primero, tenemos que los siguientes triples son ciertos:
        \begin{gather*}
            \{x=y\}\ x=x-1\ \{x+1=y\} \\
            \{x+1=y\}\ x=x+1\ \{x=y\}
        \end{gather*}
        Usando la regla de la composición, se tiene.
    \item Respecto al segundo, tenemos que los siguientes triples son ciertos:
        \begin{gather*}
            \{y=x\}\ y=y-1\ \{y+1=x\} \\
            \{y+1=x\}\ y=y+1\ \{y=x\}
        \end{gather*}
        Usando la regla de la composición, se tiene.
\end{enumerate}
Ahora, los triples son libres de interferencia por tener variables disjuntas. Podemos aplicar por tanto la regla de la composición concurrente, llegando a lo que queríamos probar:
\begin{gather*}
    \{x=y\} \\
    \texttt{cobegin} \\
    <x=x-1>;<x=x+1>;||<y=y-1>;<y=y+1>; \\
    \texttt{coend} \\
    \{x=y\}
\end{gather*}
\end{ejercicio}

\begin{ejercicio}
    Usando la regla de la conjunción, demostrar que 
    \begin{equation*}
        \{i>2\}\ i=2\ast i\ \{i>4\}
    \end{equation*}
    % // TODO: JJ. Regla de la conjunción?. Se demuestra directamente con el axioma de asignación
\end{ejercicio}

\begin{ejercicio}
    Se dan los siguientes triples de Hoare:
    \begin{gather*}
        \{j>1\}\ i=i+2;\ j=j+3;\ \{j>4\} \\
        \{i>2\}\ i=i+2;\ j=j+3;\ \{i>4\}
    \end{gather*}
    Demostrar que estos triples implican que
    \begin{equation*}
        \{j>1 \land i>2\}\ i=i+2;\ j=j+3\ \{j>4 \land i>4\}
    \end{equation*}
    ¿Qué regla se debe utilizar para la demostración?

    % // TODO: JJ. Cómo??
\end{ejercicio}

\begin{ejercicio}
    Sean $A$ y $B$ los valores iniciales de $a$ y $b$ respectivamente. Escribir un fragmento de código que tenga $\{a=A+B \land b =A-B\}$ como poscondición y demostrar que el código es correcto.\\

    En este caso, nos piden un código $C$ que cumpla:
    \begin{equation*}
        \{a=A \land b=B\}\ C\ \{a=A+B \land b =A-B\}
    \end{equation*}

    Sea $C=<a=a+b; b=a-2b>$. Buscamnos entonces demostrar los siguientes triples:
    \begin{gather*}
        \{a=A \land b=B\}\ a=a+b\ \{a=A+B \land b=B\} \\
        \{a=A+B \land b=B\}\ b=a-2b\ \{a=A+B \land b=A-B\}
    \end{gather*}

    Demostramos cada uno por separado:
    \begin{enumerate}
        \item Usando el axioma de asignación:
            \begin{align*}
                \{a=A+B \land b=B\}_{a+b}^a
                &\equiv \{a+b=A+B \land b=B\}
                \equiv\\&\equiv \{a=A\land b=B\}\ a=a+b\ \{a=A+B \land b=B\}
            \end{align*}
        \item Usando el axioma de asignación:
            \begin{align*}
                \{a=A+B &\land b=A-B\}_{a-2b}^b
                \equiv \{a=A+B \land a-2b=A-B\}
                \equiv\\& \equiv \{a=A+B \land A+B-2b=A-B\}
                \equiv\\& \equiv \{a=A+B \land b=B\}\ b=a-2b\ \{a=A+B \land b=A-B\}
            \end{align*}
    \end{enumerate}
    Usando la regla de la composición, tenemos que el código es correcto.

\end{ejercicio}

\begin{ejercicio}
    Demostrar que la siguiente sentencia tiene la poscondición ${\{x\geq 0, x^2\geq a^2\}}$.
    \verb|if a > 0 then x = a else x = -a|. Es decir, probar el triple:
    \begin{equation*}
        \{V\}\ \texttt{if a >\ 0 then x = a else x = -a}\ \{x\geq 0, x^2\geq a^2\}
    \end{equation*}~\\
    Para ello, tenemos que usar la regla del \verb|if|:
    \begin{equation*}
        \dfrac{\{P\land B\}S_1\{Q\}, \{P\land \lnot B\}S_2\{Q\}}{\{P\}\ \texttt{if B then}\ S_1\ \texttt{else}\ S_2\ \{Q\}}
    \end{equation*}
    Luego bastará con probar los triples
    \begin{gather*}
        \{a>0\} \equiv \{V \land a > 0\}\ x=a\ \{x\geq 0 \land x^2\geq a^2\} \\
        \{a\leq 0\} \equiv \{V \land a \leq 0\}\ x=-a\ \{x\geq 0 \land x^2\geq a^2\}
    \end{gather*}
    \begin{enumerate}
        \item Usando el axioma de asignación:
            \begin{equation*}
                \{a \geq 0\} \equiv \{a \geq 0 \land a^2 \geq a^2\} \equiv \{x\geq 0 \land x^2\geq a^2\}^x_{a}\ x=a\ \{x\geq 0 \land x^2\geq a^2\}
            \end{equation*}
            Como $\{a>0\}\rightarrow\{a\geq 0\}$, usamos la segunda regla de la consecuencia y tenemos el primer triple demostrado.
        \item Usando el axioma de asignación:
            \begin{equation*}
                \{a \leq 0\} \equiv \{-a \geq 0 \land a^2 \geq a^2\} \equiv \{x\geq 0 \land x^2\geq a^2\}^x_{-a}\ x=-a\ \{x\geq 0 \land x^2\geq a^2\}
            \end{equation*}
    \end{enumerate}
    Y acabamos de probar el triple que nos pedía el ejercicio.
\end{ejercicio}

\begin{ejercicio}
    El siguiente fragmento de código tiene $\{P\}\equiv \left\{sum = \frac{j(j-1)}{2}\right\}$ como precondición y poscondición. Demostrar que es verdadero:
    \begin{equation*}
        \{P\}\ sum=sum+j;\ j=j+1;\ \{P\}
    \end{equation*}

    Queremos demostrar el triple:
    \begin{equation*}
        \left\{sum=\frac{j(j-1)}{2} \right\}\ sum=sum+j;\ j=j+1;\ \left\{sum=\frac{j(j-1)}{2}\right\}
    \end{equation*}
    Para ello, será suficiente con demostrar los triples
    \begin{gather*}
        \left\{sum = \frac{j(j-1)}{2} \right\}\ sum = sum + j;\ \left\{sum = \frac{(j+1)j}{2} \right\} \\
        \left\{sum = \frac{(j+1)j}{2} \right\}\ j=j+1;\ \left\{sum = \frac{j(j-1)}{2}\right\}
    \end{gather*}
    y aplicar la regla de composición.

    \begin{enumerate}
        \item Para demostrar el primer triple, usamos el axioma de asignación:
            \begin{equation*}
                \left\{sum = \frac{(j+1)j}{2} \right\}^{sum}_{sum+j}\ sum = sum + j;\ \left\{sum = \frac{(j+1)j}{2} \right\} 
            \end{equation*}
            Usando la definición de Sustitución Textual, tenemos que:
            \begin{gather*}
                \left\{sum = \frac{(j+1)j}{2} \right\}^{sum}_{sum+j} \equiv \left\{sum+j = \frac{(j+1)j}{2} \right\} \equiv \\ 
                \equiv \left\{sum = \frac{(j+1)j}{2}-j \right\} \equiv \left\{sum = \frac{j(j-1)}{2} \right\}
            \end{gather*}
        \item Para el segundo, usamos también el axioma de asignación:
            \begin{equation*}
                \left\{sum = \frac{j(j-1)}{2}\right\}^j_{j+1} j=j+1;\ \left\{sum = \frac{j(j-1)}{2}\right\}
            \end{equation*}
            Usando de nuevo la definición de Sustitución Textual, tenemos que:
            \begin{gather*}
                \left\{sum = \frac{j(j-1)}{2}\right\}^j_{j+1} \equiv \left\{sum = \frac{(j+1)(j+1-1)}{2}\right\} \equiv \left\{sum = \frac{(j+1)j}{2}\right\}
            \end{gather*}
    \end{enumerate}
    Por lo que el triple del enunciado es cierto.

\end{ejercicio}

\begin{ejercicio}
    Demostrar que
    \begin{equation*}
        \{i\ast j + 2 \ast j + 3\ast i = 0\}\ j=j+3;\ i=i+2;\ \{i\ast j=6\}
    \end{equation*}

    Vamos buscando aplicar la regla de la composición. Para ello, y como desconocemos el estado intermedio por el que debemos pasar, usamos directamente la Sustitución Textual al final, para así obtener la precondición del segundo triple.
    \begin{equation*}
        \{i\ast j=6\}_{i+2}^i \equiv \{(i+2)\ast j=6\} \equiv \{i\ast j+2\ast j=6\} \equiv \{j\ast(i+2)=6\}
    \end{equation*}

    Usando esa precondición, el segundo libre se demuestra directamente con el axioma de asignación. Demostramos ahora el primer triple:
    \begin{gather*}
        \{i\ast j + 2 \ast j + 3\ast i = 0\}\ j=j+3;\ \{j\ast(i+2)=6\}
    \end{gather*}

    Usando la sustitución textual, tenemos que:
    \begin{align*}
        \{j\ast(i+2)=6\}_{j+3}^j &\equiv \{(j+3)\ast(i+2)=6\} \equiv
        \{j\ast(i+2)+3\ast(i+2)=6\} \equiv \\ &\equiv \{i\ast j + 2 \ast j + 3\ast i = 0\}
    \end{align*}

    Por lo que, tras usar la regla de la composición, vemos que el triple del enunciado es cierto.
\end{ejercicio}

\begin{ejercicio}
    ¿Por qué en la regla del \verb|while B|, la condición \verb|B| debe ser verdadera al comienzo del bucle?

    % // TODO: No entiendo a qué se refiere
\end{ejercicio}

\begin{ejercicio}
    Considerar una función con dos argumentos que se usa en un programa. Explicar por qué el uso de alias puede ser un problema en este caso.

    % // TODO: No entiendo a qué se refiere
\end{ejercicio}

\begin{ejercicio}
    Demostrar la corrección parcial del siguiente fragmento de programa:
    \begin{minted}{pascal}
        sum:= 0; j:= 1;
        while j <> c do begin {<> es !=}
            sum:= sum+j;
            j:= j+1;
        end
        {sum = c*(c-1)/2}
    \end{minted}

    Para ello, tenemos que hacer uso de la regla de la iteración: 
    \begin{equation*}
        \dfrac{\{I\land B\}S\{I\}}{\{I\}\ \texttt{while B do}\ S\ \texttt{end do}\ \{I\land \lnot B\}}
    \end{equation*}

    Identificando términos, sean:
    \begin{align*}
        I &\equiv sum = \frac{j(j-1)}{2} \\
        B &\equiv j \neq c \\
        S &= sum = sum+j;\ j = j+1
    \end{align*}
    Luego tendremos que probar que se cumple el triple
    \begin{equation*}
        \left\{sum=\frac{j(j-1)}{2} \land j\neq c\right\}\ sum=sum+j;j=j+1;\ \left\{sum=\frac{j(j-1)}{2}\right\}
    \end{equation*}
    Para ello, será suficiente con demostrar los triples
    \begin{gather*}
        \left\{sum = \frac{j(j-1)}{2} \land j\neq c\right\}\ sum = sum + j;\ \left\{sum = \frac{(j+1)j}{2} \land j \neq c\right\} \\
        \left\{sum = \frac{(j+1)j}{2} \land j\neq c\right\}\ j=j+1;\ \left\{sum = \frac{j(j-1)}{2}\right\}
    \end{gather*}
    y aplicar la regla de composición.

    \begin{enumerate}
        \item Para demostrar el primer triple, usamos el axioma de asignación:
            \begin{equation*}
                \left\{sum = \frac{(j+1)j}{2} \land j \neq c\right\}^{sum}_{sum+j}\ sum = sum + j;\ \left\{sum = \frac{(j+1)j}{2} \land j \neq c\right\} 
            \end{equation*}
            \begin{gather*}
                \left\{sum = \frac{(j+1)j}{2} \land j \neq c\right\}^{sum}_{sum+j} \equiv \left\{sum+j = \frac{(j+1)j}{2} \land j \neq c\right\} \equiv \\ 
                \equiv \left\{sum = \frac{(j+1)j}{2}-j \land j \neq c\right\} \equiv \left\{sum = \frac{j(j-1)}{2} \land j \neq c\right\}
            \end{gather*}
        \item Para el segundo, usamos también el axioma de asignación:
            \begin{equation*}
                \left\{sum = \frac{j(j-1)}{2}\right\}^j_{j+1} j=j+1;\ \left\{sum = \frac{j(j-1)}{2}\right\}
            \end{equation*}
            \begin{gather*}
                \left\{sum = \frac{j(j-1)}{2}\right\}^j_{j+1} \equiv \left\{sum = \frac{(j+1)(j+1-1)}{2}\right\} \equiv \left\{sum = \frac{(j+1)j}{2}\right\}
            \end{gather*}
            Además, tenemos que se tiene:
            \begin{equation*}
                \left\{sum = \frac{(j+1)j}{2} \land j \neq c\right\} \rightarrow \left\{sum = \frac{(j+1)j}{2}\right\}
            \end{equation*}
            Por tanto, usando la segunda regla de la consecuencia, tenemos que el segundo triple es cierto.
    \end{enumerate}

    Por tanto, mediante la regla de la iteración, tenemos que:
    \begin{equation*}
        \left\{sum=\frac{j(j-1)}{2}\right\}\ \texttt{while j <> c do begin sum:= sum+j; j:= j+1 end}\ \left\{sum=\frac{c(c-1)}{2}\right\}
    \end{equation*}
    Como inicialmente $j=1$ y $sum=0$, tenemos que este triple coincide con el enunciado del ejercicio.
\end{ejercicio}

\begin{ejercicio}
    Demostrar la corrección del siguiente triple:
    \begin{equation*}
        \{a[i]\geq 0\}\ a[i]=a[i]+a[j];\ \{a[i]\geq a[j]\}
    \end{equation*}

    Para ello, basta aplicar el axioma de la asignación:
    \begin{equation*}
        \{a[i]\geq a[j]\}_{a[i]+a[j]}^{a[i]}\ a[i]=a[i]+a[j];\ \{a[i]\geq a[j]\}
    \end{equation*}
    \begin{equation*}
        \{a[i]\geq a[j]\}_{a[i]+a[j]}^{a[i]}\ \equiv \{a[i]+a[j]\geq a[j]\} \equiv \{a[i] \geq 0\}
    \end{equation*}
\end{ejercicio}

\begin{ejercicio}
    Verificar el siguiente segmento de programa:
    \begin{gather*}
        \{n\geq 0\}\\ 
        i = 1;\\
        \texttt{while\ } i\leq n \texttt{\ do\ begin}\\
        a[i] = b[i]; \\
        i=i+1;\\
        \texttt{end} \\
        \left\{\bigwedge_{i=1}^n (a[i]=b[i])\right\}
    \end{gather*}

    % // TODO: Arturo se ha quedado por este ejercicio
\end{ejercicio}

\begin{ejercicio}
    El siguiente fragmento de programa calcula $\displaystyle\sum_{i=1}^{n} i!$. Demostrar que es correcto.
    \begin{minted}{pascal}
        i = 1; sum = 0; f = 1;
        while i <> n+1 do begin
          sum = sum + f;
          i = i + 1;
          f = f * 1;
        end
    \end{minted}
\end{ejercicio}

\begin{ejercicio}
    Hallar la precondición $\{P\}$ que hace que el siguiente triple sea correcto:
    \begin{equation*}
        \{P\}\ a[i]=2\ast b;\ \{j\leq i ~\land~ k < i ~\land~ a[i] + a[j-1]+a[k]>b\}
    \end{equation*}
    Para ello, basta aplicar el axioma de asignación:
    \begin{gather*}
        \{j\leq i~\land~ k < i ~\land~ a[i] + a[j-1]+a[k]>b\}_{2\cdot b}^{a[i]} a[i]=2\ast b;\\ \{j\leq i ~\land~ k < i ~\land~ a[i] + a[j-1]+a[k]>b\}
    \end{gather*}

    Usando la definición de Sustitución Textual, tenemos que:
    \begin{gather*}
        \{j\leq i~\land~ k < i ~\land~ a[i] + a[j-1]+a[k]>b\}_{2\cdot b}^{a[i]} \equiv \\ \equiv \{j\leq i ~\land~ k < i ~\land~ 2\cdot b + a[j-1]+a[k]>b\}
        \ \equiv\\ \equiv \{j\leq i ~\land~ k < i ~\land~ b + a[j-1]+a[k]>0\}
    \end{gather*}
    Luego estamos buscando la precondición:
    \begin{equation*}
        \{P\} \equiv \{j\leq i ~\land~ k < i ~\land ~b + a[j-1]+a[k]>0\}
    \end{equation*}
\end{ejercicio}

\begin{ejercicio}
    Demostrar que para $n>0$ el siguiente fragmento de programa termina.
    \begin{minted}{pascal}
        i = 1; f = 1;
        while i <> n do begin
          i = i + 1;
          f = f * r;
        end
    \end{minted}
    % // TODO: JJ. Cómo demostrar que termina?
\end{ejercicio}

\begin{ejercicio}
    Hallar la precondición de la terna:
    \begin{equation*}
        \{P\}\ a[i]=b;\ \{a[j] = 2\ast a[i]\}
    \end{equation*}
    Para ello, simplemente aplicamos el axioma de asignación:
    \begin{equation*}
        \{a[j]=2\ast b\} \equiv \{a[j]=2\ast a[i]\}_{b}^{a[i]}\ a[i]=b;\ \{a[j] = 2\ast a[i]\}
    \end{equation*}
    % // TODO: JJ. Tenemos que considerar i=j?
\end{ejercicio}

\begin{ejercicio}
    Para cada uno de los siguientes fragmentos de código, obtener la poscondición apropiada:
    \begin{enumerate}
        \item $\{i<10\}\ i=2\ast i+1;$.\\
            La poscondición es $\{i<21\}$:
            \begin{equation*}
                \{i<10\}\ i=2\ast i+1; \{i<21\}
            \end{equation*}
            que puede demostrarse aplicando el axioma de asignación.
        \item $\{i>0\}\ i=i-1;$.\\
            La poscondición es $\{i>-1\}$:
            \begin{equation*}
                \{i>0\}\ i=i-1; \{i>-1\}
            \end{equation*}
            que puede demostrarse aplicando el axioma de asignación.
        \item $\{i>j\}\ i=i+1;\ j=j+1;$.\\
            La poscondición es $\{i>j\}$:
            \begin{gather*}
                \{i>j\}\ i=i+1;\ \{i>j+1\} \\
                \{i>j+1\}\ j=j+1;\ \{i>j\}
            \end{gather*}
            Ambos pueden demostrarse aplicando el axioma de asignación y finalmente tenemos que:
            \begin{equation*}
                \{i>j\}\ i=i+1;\ j=j+1;\ \{i>j\}
            \end{equation*}
            aplicando la regla de composición.
        \item $\{V\}\ i=3;\ j=2\ast i$.\\
            La poscondición es $\{i=3 \land j=6\}$:
            \begin{gather*}
                \{V\}\ i=3;\ \{i=3\} \\
                \{i=3\}\ j=2\ast i;\ \{i=3 \land j = 6\}
            \end{gather*}
    \end{enumerate}
\end{ejercicio}

\begin{ejercicio}
    Para cada uno de los siguientes fragmentos de código, obtener las precondiciones apropiadas.
    \begin{enumerate}
        \item $i=3\ast k;\ \{i>6\}$.\\
            Aplicando el axioma de asignación:
            \begin{equation*}
                \{k>2\} \equiv \{3\cdot k>6\} \equiv \{i>6\}^i_{3\cdot k}\ i=3\ast k;\ \{i>6\}
            \end{equation*}
            obtenemos que la precondición es $\{k>2\}$.
        \item $a=b\ast c;\ \{a=1\}$.\\
            Aplicando el axioma de asignación:
            \begin{equation*}
                \{b = c^{-1}\} \equiv \{b\cdot c = 1\} \equiv \{a=1\}^a_{b\cdot c}\ a=b\ast c;\ \{a=1\}
            \end{equation*}
        La precondición es $\{b=c^{-1}\}$.
        \item $b=c-2;\ a=a/b;$\footnote{Esto es lo que pone en la relación, falta la poscondición para poder hacerse.}
    \end{enumerate}
\end{ejercicio}

\begin{ejercicio}
    Verificar el siguiente código. Indicar todas las reglas usadas.
    \begin{equation*}
        \{y>0\}\ xa=x+y;\ xb=x-y;
    \end{equation*}
    % // TODO: JJ. No hay poscondición??
\end{ejercicio}

\begin{ejercicio}
    Verificar el siguiente código, indicando todas las reglas usadas.
    \begin{equation*}
        \{V\}\ if\ x < 0\ then\ x = -x\ \{x\geq 0\}
    \end{equation*}
    Para comenzar, probamos que $\{x<0\}\ x=-x;\ \{x\geq 0\}$ usando el axioma de asignación:
    \begin{equation*}
        \{x<0\} \equiv \{-x\geq 0\} \equiv \{x\geq 0\}^x_{-x}\ x=-x;\ \{x\geq 0\}
    \end{equation*}
    Posteriormente, como sabemos que $\{V \land (x<0)\}\equiv \{x<0\}$ y que ${\{x\geq 0\}\ null\ \{x\geq 0\}}$ por el axioma de la sentencia nula, podemos aplicar la regla del \texttt{if}:
    \begin{equation*}
        \dfrac{\{V \land (x<0)\}\ x=-x;\ \{x\geq 0\}, \{V \land x\geq 0\}\ null\ \{x\geq 0\}}{\{V\}\ \texttt{if\ }x<0 \texttt{\ then\ }x=-x \texttt{\ else\ null\ } \{x\geq 0\}}
    \end{equation*}
\end{ejercicio}

\begin{ejercicio}\label{ej_rel1_44}
    Verificar el siguiente segmento de programa:
    \begin{gather*}
        max = a[1]; i = 1; \\
        \texttt{while\ } i <> n+1 \texttt{\ do\ begin} \\
        \texttt{if\ } a[i] \geq max \texttt{\ then\ } max = a[i];\\
        i=i+1; \\
        \texttt{end} \\
        \left\{\bigwedge_{i=1}^n (max \geq a[i])\right\}
    \end{gather*}
\end{ejercicio}

\begin{ejercicio}
    Demostrar la corrección parcial del siguiente código:
    \begin{gather*}
        max = a[1]; i = 1; \\
        \texttt{while\ } i < n \texttt{\ do\ begin} \\
        i = i +1; \\
        \texttt{if\ } a[i] \geq max \texttt{\ then\ } max = a[i]; \\
        \texttt{end} \\
        \left\{\bigwedge_{i=1}^n (max \geq a[i]), \bigvee_{j=1}^n (max = a[i]) \right\}
    \end{gather*}
\end{ejercicio}

\begin{ejercicio}
    Demostrar la corrección parcial del siguiente código:
    \begin{gather*}
        i = 0; j= n; \\
        \texttt{while\ } i < n \texttt{\ do\ begin} \\
        i = i + 1; \\
        j = j - 1; \\
        a[i] = b[j] \\
        \texttt{end} \\
        \left\{\bigwedge_{i=1}^n(a[i] = b[n-i])\right\}
    \end{gather*}
\end{ejercicio}

\begin{ejercicio}
    Demostrar la corrección parcial del siguiente código:
    \begin{minted}{pascal}
        i = 0;
        s = 0;
        while i <= n do begin
          s = s + a[i];
          a[i] = s;
          i = i + 1;
        end
    \end{minted}
\end{ejercicio}

\begin{ejercicio}
    Dados $n\geq 0$, $i\leq n$, demostrar que el siguiente segmento de progama evalúa
    \begin{equation*}
        \dfrac{n!}{i!(n-i)!}
    \end{equation*}
    \begin{minted}{pascal}
        k = 0; fact = 1;
        while k <> n do begin
          k = k + 1;
          fact = fact * k;
          if k <= i then afact = fact;
          if k <= n-i then bfact = fact;
        end
        bcof = fact/(afact*bfact);
    \end{minted}
\end{ejercicio}

\begin{ejercicio}
    Demostrar la terminación del fragmento de programa dado en el problema~\ref{ej_rel1_44} ¿Qué condición se debe imponer para realizar la demostración?
\end{ejercicio}

\end{document}
